% =====================================================================
% Chapter 2 -- The flow of formalisations
% Sources: FACTS.md; papers/ar_dipolar_exciton, ar_yukawa_crystal, ar_keldysh_tmd,
% prog_cm_bernstein_library, prog_sandier_serfaty_audit, p2_sandier_serfaty_lean, p2_3d_feasibility;
% Lean: formal_cert/Autoresearch/{DipolarExciton*,Yukawa*,KeldyshCM*}.lean,
% formal_cert/Programmes/CompletelyMonotone/, formal_cert/Programmes/SandierSerfaty/SignedEndpoints.lean.
% The Lean excerpts are VERBATIM; the line number refers to the source file.
% =====================================================================

% Chapter-local listings style: the "literate" list of main.tex, copied as is,
% extended by a few symbols that occur in the excerpts of this chapter (β, ∫, ⁻, ᵐ, ι, ᵢ).
\lstdefinestyle{chdeux}{language=Lean,
  literate={⁻}{{$^{-}$}}1 {⁻¹}{{$^{-1}$}}1 {ℝ}{{$\mathbb{R}$}}1 {ℤ}{{$\mathbb{Z}$}}1 {ℕ}{{$\mathbb{N}$}}1 {≤}{{$\le$}}1 {≥}{{$\ge$}}1 {≠}{{$\ne$}}1
  {→}{{$\to$}}1 {↔}{{$\leftrightarrow$}}1 {∀}{{$\forall$}}1 {∃}{{$\exists$}}1 {∧}{{$\wedge$}}1 {∨}{{$\vee$}}1
  {¬}{{$\neg$}}1 {⊤}{{$\top$}}1 {∈}{{$\in$}}1 {∉}{{$\notin$}}1 {∩}{{$\cap$}}1 {√}{{$\surd$}}1
  {π}{{$\pi$}}1 {ρ}{{$\rho$}}1 {κ}{{$\kappa$}}1 {α}{{$\alpha$}}1 {δ}{{$\delta$}}1 {ε}{{$\varepsilon$}}1 {σ}{{$\sigma$}}1
  {•}{{$\bullet$}}1 {∑}{{$\sum$}}1 {²}{{$^2$}}1 {₀}{{$_0$}}1 {₁}{{$_1$}}1 {₂}{{$_2$}}1 {‖}{{$\|$}}1 {⟨}{{$\langle$}}1 {⟩}{{$\rangle$}}1
  {𝓝}{{$\mathcal{N}$}}1 {∞}{{$\infty$}}1 {×}{{$\times$}}1 {·}{{$\cdot$}}1 {∘}{{$\circ$}}1 {∖}{{$\setminus$}}1
  {ψ}{{$\psi$}}1 {φ}{{$\varphi$}}1 {λ}{{$\lambda$}}1 {μ}{{$\mu$}}1 {θ}{{$\theta$}}1 {Λ}{{$\Lambda$}}1 {∂}{{$\partial$}}1
  {—}{{---}}1 {–}{{--}}1 {é}{{\'e}}1 {ö}{{\"o}}1 {’}{{'}}1
  {β}{{$\beta$}}1 {∫}{{$\int$}}1 {ᵐ}{{$^{\mathrm{m}}$}}1 {ι}{{$\iota$}}1 {ᵢ}{{$_i$}}1 {△}{{$\triangle$}}1
  {⊥}{{$\bot$}}1 {χ}{{$\chi$}}1}

\chapter{The flow of formalisations: from universal optimality to physical potentials}\label{ch2:chapitre}
\chaptermark{The flow of formalisations}

\begin{center}
\begin{minipage}{0.9\linewidth}\small\itshape
In Chapter~1, a theorem from elsewhere --- the triangular lattice minimises the energy of every potential that is a
completely monotone function of the squared distance --- was stated, read and checked by the Lean kernel. This chapter
tells what happened next: how, following a method fixed in advance, that theorem was applied to three interactions from
materials physics; how a small library turned these applications into a few lines each; and how the same method, turned
towards a neighbouring and much harder problem (the Sandier--Serfaty conjecture), served above all to state precisely
what is \emph{not} proved.
\end{minipage}
\end{center}

\begin{warnbox}[: read before the whole chapter]
\begin{itemize}[leftmargin=1.2em,itemsep=2pt]
\item Everything that rests on openai/math is \emph{conditional} on the upstream theorems, which are AI-generated and
\emph{not peer reviewed}.
\item Every result in this chapter was produced by AI agents (Claude) under the author's direction, then audited by a
\emph{separate} AI agent. \textbf{No human has audited these results}, neither the faithfulness of the Lean statements
nor the ``human'' proofs written here.
\item A \tierA{} concerns the Lean statement, never its physical reading.
\item The Sandier--Serfaty conjecture is \textbf{not proved} in this book (Section~\ref{ch2:sec:ss}).
\end{itemize}
\end{warnbox}

% ---------------------------------------------------------------------
\section{Overview: a flow, not a catalogue}\label{ch2:sec:flux}

The starting point (Chapter~1) is the upstream theorem \lean{universal_energy_minimum} of openai/math: for every
\emph{admissible} potential $g$, the triangular lattice $\Lat$ of covolume one has energy less than or equal to that of
every locally finite configuration of density one \tierA{}, conditionally on the upstream~\cite{oai}. In the human
literature, universal optimality is proved in dimensions 8 and 24~\cite{ckmrv} \tierL, and Petrache and Serfaty wrote in
2019 that the planar case remained open~\cite{ch2ps} \tierL: the upstream theorem is therefore an AI claim on an open
problem, accepted here only as a Lean statement checked by the kernel from the upstream definitions. We recall the two
statements, quoted verbatim.

\begin{leanbox}[--- openai/math, \lean{OAI/Analysis/Triangular/Energy/Basic.lean:26-29}]
\begin{lstlisting}[style=chdeux]
def AdmissiblePotential (g : ℝ → ℝ) : Prop :=
  ContDiffOn ℝ (⊤ : ℕ∞) g (Set.Ioi 0) ∧
  (∀ t : ℝ, 0 < t → 0 ≤ g t) ∧
  ∀ (r : ℕ) (t : ℝ), 0 < t → 0 ≤ (-1 : ℝ) ^ r * iteratedDeriv r g t
\end{lstlisting}
\end{leanbox}
\begin{leanbox}[--- openai/math, \lean{OAI/Analysis/Triangular/Energy/Universal.lean:9-11} (proof omitted)]
\begin{lstlisting}[style=chdeux]
theorem universal_energy_minimum (g : ℝ → ℝ) (C : Set Plane)
    (hg : AdmissiblePotential g) (_hC : LocallyFinite C) (hd : DensityOne C) :
    latticeEnergy g ≤ energy g C ∧ latticeEnergy g = energy g A := by
\end{lstlisting}
\end{leanbox}
\noindent (These two excerpts are quoted from the pinned copy of openai/math; the working copy
\lean{results/openai_math/hypotheses/night_2026-10-08/TRI_D3/src/} in the repository is identical to it for these two files.) The variable of $g$ is the \emph{square} $t=r^2$ of the
distance; ``admissible'' means: $C^\infty$ on $(0,+\infty)$, nonnegative, and $(-1)^r g^{(r)}(t)\ge0$ for every $r$,
that is, \emph{completely monotone} (CM) on $(0,+\infty)$. The energies take values in $[0,+\infty]$.

The flow of this chapter is as follows (Figure~\ref{ch2:fig:flux}).
\begin{enumerate}[itemsep=1pt]
\item \textbf{Upstream} (Chapter~1): universal optimality in dimension two, \tierA{} conditional, and our corollaries
D1--D4 (Riesz, every density, CKMRV sense, open or closed disks), published in textbook~v1
(doi\string:10.5281/zenodo.23295142)~\cite{v1}.
\item \textbf{Three physical potentials} (Section~\ref{ch2:sec:pot}): dipolar excitons, planar Yukawa, Rytova--Keldysh.
For each: a proof of complete monotonicity in $r^2$, the Lean statement, the consequence.
\item \textbf{The generalisation} (Section~\ref{ch2:sec:lib}): a CM/Bernstein library whose subordination theorem turns
these three results into corollaries of two or three lines.
\item \textbf{A neighbouring problem} (Section~\ref{ch2:sec:ss}): the Sandier--Serfaty conjecture, where the same method
produces a reproduction, an audit, a Lean statement, a \emph{partial} formalisation and a negative result in dimension
three --- but no proof.
\end{enumerate}

\begin{figure}[ht]\centering
\begin{tikzpicture}[x=1cm,y=1cm,font=\small,
  bx/.style={draw,rounded corners=2pt,align=center,inner sep=4pt,fill=#1!6},
  ar/.style={->,>=stealth,thick,gray!70!black}]
\node[bx=tierL] (up) at (0,0) {openai/math\\ \lean{universal_energy_minimum}\\ (AI, not peer reviewed)};
\node[bx=tierA] (d) at (4.6,0) {D1--D4\\ (Ch.~1)};
\node[bx=tierA] (h1) at (0,-2.2) {H1 dipolar\\ $2(t^{-1/2}-(t+d^2)^{-1/2})$};
\node[bx=tierA] (h2) at (4.6,-2.2) {H2 Yukawa\\ $e^{-\kappa\sqrt t}/\sqrt t$};
\node[bx=tierA] (h3) at (9.2,-2.2) {H3 Keldysh\\ Laplace mixture};
\node[bx=tierA] (p3) at (4.6,-4.4) {P3: CM/Bernstein library\\ subordination $f\circ\psi$};
\node[bx=tierC] (ss) at (9.2,0) {Sandier--Serfaty\\ P1, A1, B3\\ (not proved)};
\draw[ar] (up) -- (d);
\draw[ar] (up) -- (h1);
\draw[ar] (up) -- (h2);
\draw[ar] (up.south east) -- (h3.north west);
\draw[ar] (h1) -- (p3);
\draw[ar] (h2) -- (p3);
\draw[ar] (h3) -- (p3);
\draw[ar,dashed] (d) -- node[above,font=\scriptsize]{same method} (ss);
\end{tikzpicture}
\caption{The flow of formalisations in this chapter. Green: \tierA{} Lean statements (conditional on the upstream);
blue: cited upstream; orange: work whose mathematical status remains \tierC{} or lower.}\label{ch2:fig:flux}
\end{figure}

% ---------------------------------------------------------------------
\section{The method: preregister, falsify, let the kernel judge}\label{ch2:sec:methode}

Before the first computation, one must say how one will know that one was wrong. The method followed in all the work of
this chapter comes down to four rules. It is tooled by Elenchus~\cite{elenchus}, a claims register (``ledger'') equipped
with an automatic gate.

\subsection{Rule 1: preregister}\label{ch2:ssec:prereg}
One writes down, \emph{before} seeing the results, what one is looking for, what one will report, and what would make the
project fail. Two real examples.
\begin{itemize}
\item For the dipolar excitons, the preregistration \texttt{DIP-TRI-1} was committed (commit
\texttt{4feb027}) \emph{before} the solver code (commit \texttt{17ea6be}) in rusty-SUNDIALS~\cite{ch2dip}. It fixed
the quantity to be reported (a truncated image sum plus the uniform tail). The paper reports exact Ewald energies
instead, and says so explicitly as a \emph{deviation}: the truncated sum left a non-monotone residual.
\item For the formalisation of the Sandier--Serfaty certificate (A1), the prior probabilities were written in
\texttt{preregistration.json}: $0.6$ that a few pilot stages succeed, $0.2$ that all of them succeed, $0.35$
for one complete lemma of the reduction, $0.02$ for the whole; and a \emph{stopping criterion}: no announcement of a formal
proof as long as all stages and the whole reduction are not checked by the kernel~\cite{ch2ssl}. The criterion
applies (Section~\ref{ch2:ssec:a1}).
\end{itemize}

\subsection{Rule 2: a falsification gate, with positive and negative controls}\label{ch2:ssec:falsif}
Every experiment, every theorem, every program comes with a \emph{positive control} (a case where the answer must be
``yes'') and a \emph{negative control} (a case where it must be ``no''), run before any number is reported. A gate that
accepts everything checks nothing. In this chapter, every admissibility statement has its counterexample checked by the
kernel: $t^{-1/2}-2(t+d^2)^{-1/2}$ for the dipoles, $e^{-\kappa\sqrt t}\cos t$ and $g_1-2g_2$ for
Yukawa, $e^{-t}-2e^{-2t}$ for the Laplace mixtures, $\cos$ and a ``junk'' composition for the library.
For the Sandier--Serfaty certificate, three deliberately weakened versions of the program were run; Section~\ref{ch2:ssec:p1}
says honestly that one of them \emph{passes}.

\subsection{Rule 3: the kernel is the judge --- but not its exit code}\label{ch2:ssec:noyau}
A Lean file that contains \lean{sorry}, or a planted axiom, \emph{compiles with exit code 0}. The exit code is therefore
not the gate. Here is the planted-witness file of the CM library, quoted verbatim.
\begin{leanbox}[--- \lean{formal_cert/Programmes/CompletelyMonotone/Controls.lean:11-18}]
\begin{lstlisting}[style=chdeux]
theorem ProgCM.ctl_sorry : ProgCM.IsCompletelyMonotoneOn Real.cos := sorry

axiom ProgCM.ctl_cheat : False

theorem ProgCM.ctl_axiom : ProgCM.IsCompletelyMonotoneOn Real.cos := ProgCM.ctl_cheat.elim

#print axioms ProgCM.ctl_sorry
#print axioms ProgCM.ctl_axiom
\end{lstlisting}
\end{leanbox}
\noindent Both ``proofs'' of the false statement ``$\cos$ is CM'' compile with exit code~0; only the command
\texttt{\#print axioms} unmasks them (\lean{sorryAx}, resp.\ \lean{ProgCM.ctl_cheat})~\cite{p3}. Hence the definition of
\tierA{} (fact sheet \texttt{FACTS.md}): the list of axioms is contained in $\{$\lean{propext}, \lean{Classical.choice},
\lean{Quot.sound}$\}$, plus a kernel replay (\texttt{leanchecker}), plus a textual search for
\lean{skipKernelTC}, \lean{addDecl}, \lean{native_decide}, \lean{implemented_by}.

\begin{remark}[Why \texttt{\#print axioms} is not enough]\label{ch2:rem:skip}
The A2 audit of openai's Navier--Stokes corpus showed that a declaration forged with the option
\lean{debug.skipKernelTC} can display an impeccable axiom list. A full kernel replay closes this
backdoor. In the interest of honesty: the source papers of these sections did not report a \texttt{leanchecker} replay. This replay was
carried out for this book on 2026-10-26 (record \texttt{results/kernel\_replay\_2026-10-26/}), on the modules
cited here: exit code~0 for each of the dipolar modules (\lean{DipolarExciton}, \lean{DipolarExcitonMultipolar},
\lean{DipolarExcitonStep1}), the Yukawa modules (\lean{Yukawa}, \lean{YukawaMix}), the Keldysh module (\lean{KeldyshCM}) and the
CM/Bernstein library (\lean{Basic}, \lean{Bernstein}, \lean{Laplace}, \lean{Corollaries}, \lean{Applications},
\lean{Examples}, \lean{Limits}), and the textual search finds no forbidden token outside the planted-control
files. The replay covers the declarations of these modules; Mathlib and the Lean core are loaded, not replayed.
The modules of Section~\ref{ch2:sec:ss} (Sandier--Serfaty) are in the replay plan but had not been
replayed when this book was assembled: their statements carry the label \tierAs{}.
\end{remark}

\subsection{Rule 4: the producer is not the verifier}\label{ch2:ssec:prodverif}
The agent that writes a theorem is never the one that declares it verified. A second AI agent recompiles the files
\emph{from source}, in its own directory, plants its own controls, reads the statements adversarially, and returns a
verdict (for instance \texttt{ACCEPT\_WITH\_CORRECTIONS}). The Elenchus register assigns each claim a tier --- A (kernel),
L (literature read and pinned), C (human argument not machine-checked), X (floating-point exploration, which never
supports a claim) --- and the gate checks that no row is ranked above what it rests on, transitively. It raises a flag
\texttt{LEDGER\_UNAUDITED\_TIER\_A} for every Tier~A row for which no human has certified that the statement means what
its gloss says. These flags remain open in all the work of this chapter: an AI audit does not close them.

\begin{warnbox}[: what a Tier A certifies]
The kernel certifies that \emph{a Lean statement} is proved from the Lean definitions. It certifies neither that these
definitions faithfully translate the physics, nor that the upstream (openai/math) is correct in the human sense: it
certifies that, \emph{if} one accepts the upstream definitions and theorems, the conclusion follows.
\end{warnbox}

% ---------------------------------------------------------------------
\section{Three physical potentials, three proofs}\label{ch2:sec:pot}

The three subsections follow the same plan: (i) the physical model \tierC{}; (ii) a step-by-step proof that the
potential, viewed as a function of $t=r^2$, is completely monotone \tierC{}, alongside the verified Lean statement
\tierA{}; (iii) the consequence: optimality of the triangular lattice (conditional), and finiteness of the energy when it
is proved; (iv) a negative control.

\begin{definition}[complete monotonicity]\label{ch2:def:cm}
A function $f\colon(0,+\infty)\to\R$ is \emph{completely monotone} (CM) if it is of class $C^\infty$ and, for every
$n\in\N$ and every $t>0$, $(-1)^nf^{(n)}(t)\ge0$. In other words: $f\ge0$, $f'\le0$, $f''\ge0$, $f'''\le0$, and so on.
\end{definition}

\begin{example}
$t\mapsto e^{-at}$ ($a\ge0$) is CM since $(-1)^n\frac{d^n}{dt^n}e^{-at}=a^ne^{-at}\ge0$. The function $\cos$ is not:
$\cos\pi=-1<0$. Nor is $\sqrt t$: it is increasing.
\end{example}

\begin{remark}[why the variable $t=r^2$?]
The upstream theorem is about $g(\|x-y\|^2)$. Being CM in $r^2$ is more demanding than being CM in $r$: $e^{-r}$ is CM in $r$,
and we shall see that $e^{-\sqrt t}$ is also CM in $t$, but this requires an argument (Section~\ref{ch2:ssec:yuk}).
\end{remark}

% .....................................................................
\subsection{Dipolar excitons in a bilayer}\label{ch2:ssec:dip}

\begin{physbox}[: the model \tierC]
In a heterobilayer of transition-metal dichalcogenides (for instance MoSe$_2$/WSe$_2$), an ``indirect'' exciton is an
electron in one layer, at height $d$, bound to a hole in the other layer, at height $0$, at the same in-plane position.
Two excitons at in-plane distance $r$ form four Coulomb pairs: electron--electron and hole--hole at distance $r$
(repulsion $+1/r$ each), and two cross pairs at distance $\sqrt{r^2+d^2}$ (attraction $-1/\sqrt{r^2+d^2}$ each). The
internal binding energy, a constant, is neglected. In units of $e^2/(4\pi\varepsilon)$:
\[V(r)=\frac2r-\frac2{\sqrt{r^2+d^2}},\qquad g_d(t)=2\bigl(t^{-1/2}-(t+d^2)^{-1/2}\bigr),\quad V(r)=g_d(r^2).\]
At large distance $V(r)=d^2/r^3+O(r^{-5})$: the dipole--dipole law. This pair counting is \tierC{}; the algebraic
identity $g_d(r^2)=2(1/r-1/\sqrt{r^2+d^2})$ is checked by the kernel (\lean{g_inplane}, \tierA{}: the module \texttt{DipolarExcitonStep1.lean} was verified only by its producer in the source paper~\cite{ch2dip}; it has since been replayed by \lean{leanchecker}, exit code~0).
\end{physbox}

\begin{leanbox}[--- \lean{formal_cert/Autoresearch/DipolarExciton.lean:22}]
\begin{lstlisting}[style=chdeux]
noncomputable def g (d : ℝ) : ℝ → ℝ := fun t => 2 * (t ^ (-(1 / 2 : ℝ)) - (t + d ^ 2) ^ (-(1 / 2 : ℝ)))
\end{lstlisting}
\end{leanbox}

\begin{theorem}[H1, admissibility of the dipolar potential]\label{ch2:thm:dip}
For every real $d$, $g_d$ is completely monotone on $(0,+\infty)$. More generally, for $s>0$ and $c\ge0$,
$\mathrm{shiftDiff}_{s,c}(t)=t^{-s/2}-(t+c)^{-s/2}$ is CM. \tierA{} (Lean statement); proof below \tierC.
\end{theorem}

The proof rests on \emph{Pochhammer symbols}. For $a\in\R$ and $r\in\N$, write
\[\mathrm{desc}_r(a)=a(a-1)\cdots(a-r+1),\qquad \mathrm{asc}_r(a)=a(a+1)\cdots(a+r-1),\]
with $\mathrm{desc}_0=\mathrm{asc}_0=1$ (falling and rising factorials; Mathlib: \lean{descPochhammer},
\lean{ascPochhammer}).

\begin{lemma}[derivatives of a shifted power]\label{ch2:lem:pochder}
For $a\in\R$, $c\ge0$, $t>0$ and $r\in\N$:
$\displaystyle\frac{d^r}{dt^r}(t+c)^a=\mathrm{desc}_r(a)\,(t+c)^{a-r}$.
\end{lemma}
\begin{humanproof}
Induction on $r$. For $r=0$ it is the identity. If the formula holds at order $r$, differentiate $(t+c)^{a-r}$, which is
legitimate since $t+c>0$: one obtains $(a-r)(t+c)^{a-r-1}$, and $\mathrm{desc}_r(a)(a-r)=\mathrm{desc}_{r+1}(a)$.
\end{humanproof}
\begin{leanbox}[--- \lean{DipolarExciton.lean:32-33} and \lean{:36-37} (proofs omitted)]
\begin{lstlisting}[style=chdeux]
lemma iteratedDeriv_rpow (a : ℝ) (r : ℕ) (t : ℝ) :
    iteratedDeriv r (fun x : ℝ => x ^ a) t = (descPochhammer ℝ r).eval a * t ^ (a - r) := by
\end{lstlisting}
\begin{lstlisting}[style=chdeux]
lemma iteratedDeriv_shift (a c : ℝ) (r : ℕ) (t : ℝ) :
    iteratedDeriv r (fun x : ℝ => (x + c) ^ a) t = (descPochhammer ℝ r).eval a * (t + c) ^ (a - r) := by
\end{lstlisting}
\end{leanbox}

By linearity of differentiation, for $t>0$:
\begin{equation}\label{ch2:eq:gder}
g_d^{(r)}(t)=2\,\mathrm{desc}_r\!\left(-\tfrac12\right)\Bigl(t^{-1/2-r}-(t+d^2)^{-1/2-r}\Bigr).
\end{equation}
\begin{leanbox}[--- \lean{DipolarExciton.lean:41-44} (proof omitted)]
\begin{lstlisting}[style=chdeux]
lemma iteratedDeriv_g (d : ℝ) (r : ℕ) {t : ℝ} (ht : 0 < t) :
    iteratedDeriv r (g d) t =
      2 * ((descPochhammer ℝ r).eval (-(1 / 2 : ℝ)) *
        (t ^ (-(1 / 2 : ℝ) - r) - (t + d ^ 2) ^ (-(1 / 2 : ℝ) - r))) := by
\end{lstlisting}
\end{leanbox}

\begin{lemma}[the sign]\label{ch2:lem:sign}
$(-1)^r\,\mathrm{desc}_r(-\tfrac12)=\mathrm{asc}_r(\tfrac12)=\tfrac12\cdot\tfrac32\cdots\left(r-\tfrac12\right)>0$.
\end{lemma}
\begin{humanproof}
$(-1)^r\mathrm{desc}_r(-\tfrac12)=\prod_{k=0}^{r-1}(-1)\left(-\tfrac12-k\right)=\prod_{k=0}^{r-1}\left(\tfrac12+k\right)$,
a product of strictly positive numbers.
\end{humanproof}
\begin{leanbox}[--- \lean{DipolarExciton.lean:56-57} (proof omitted)]
\begin{lstlisting}[style=chdeux]
lemma sign_desc (r : ℕ) :
    (-1 : ℝ) ^ r * (descPochhammer ℝ r).eval (-(1 / 2 : ℝ)) = (ascPochhammer ℝ r).eval (1 / 2 : ℝ) := by
\end{lstlisting}
\end{leanbox}

\begin{lemma}[the bracket]\label{ch2:lem:crochet}
For $c\ge0$, $e\le0$ and $t>0$: $t^e-(t+c)^e\ge0$.
\end{lemma}
\begin{humanproof}
For $e\le0$, $x\mapsto x^e$ is nonincreasing on $(0,+\infty)$; and $t\le t+c$.
\end{humanproof}
\begin{leanbox}[--- \lean{DipolarExciton.lean:62-63} (proof omitted)]
\begin{lstlisting}[style=chdeux]
lemma diff_nonneg (c e : ℝ) (hc : 0 ≤ c) (he : e ≤ 0) {t : ℝ} (ht : 0 < t) :
    0 ≤ t ^ e - (t + c) ^ e := by
\end{lstlisting}
\end{leanbox}

\begin{humanproof}[Proof of Theorem~\ref{ch2:thm:dip}]
\emph{Regularity.} $t\mapsto t^{-1/2}$ and $t\mapsto(t+d^2)^{-1/2}$ are $C^\infty$ wherever their base is strictly
positive, hence on $(0,+\infty)$.
\emph{Signs.} By~\eqref{ch2:eq:gder} and Lemma~\ref{ch2:lem:sign},
\[(-1)^rg_d^{(r)}(t)=2\,\mathrm{asc}_r\!\left(\tfrac12\right)\Bigl(t^{-1/2-r}-(t+d^2)^{-1/2-r}\Bigr),\]
the product of a positive number and a nonnegative bracket (Lemma~\ref{ch2:lem:crochet} with $e=-\tfrac12-r\le0$ and
$c=d^2\ge0$). The case $r=0$ gives $g_d\ge0$. Replacing $\tfrac12$ by $\tfrac s2$ ($s>0$) and $d^2$ by $c\ge0$ changes
nothing: this is the statement for $\mathrm{shiftDiff}_{s,c}$. Finally $g_d=2\,\mathrm{shiftDiff}_{1,d^2}$ and a positive
multiple of a CM function is CM, which gives the result for \emph{every} real $d$.
\end{humanproof}

\begin{leanbox}[--- \lean{DipolarExciton.lean:68} and \lean{DipolarExcitonMultipolar.lean:37, 59, 102} (proofs omitted)]
\begin{lstlisting}[style=chdeux]
theorem dipolar_admissible {d : ℝ} (_hd : 0 < d) : AdmissiblePotential (g d) := by
\end{lstlisting}
\begin{lstlisting}[style=chdeux]
noncomputable def shiftDiff (s c : ℝ) : ℝ → ℝ := fun t => t ^ (-(s / 2)) - (t + c) ^ (-(s / 2))
\end{lstlisting}
\begin{lstlisting}[style=chdeux]
theorem shiftDiff_admissible {s c : ℝ} (hs : 0 < s) (hc : 0 ≤ c) : AdmissiblePotential (shiftDiff s c) := by
\end{lstlisting}
\begin{lstlisting}[style=chdeux]
theorem dipolar_admissible_all (d : ℝ) : AdmissiblePotential (AutoresearchDipolar.g d) := by
\end{lstlisting}
\end{leanbox}

\begin{remark}[Bernstein reading \tierC]
Equivalently, $t^{-1/2}-(t+c)^{-1/2}=\pi^{-1/2}\int_0^\infty e^{-ut}\,(1-e^{-uc})\,u^{-1/2}\,du$ is a positive mixture of
exponentials, which makes complete monotonicity evident~\cite{ch2dip}. This identity is not checked in Lean; the
Pochhammer proof above avoids it.
\end{remark}

\paragraph{Consequence 1: optimality at density one \tierA, conditional.} We apply the upstream theorem to $g_d$.
The Lean proof is one line: the first component of \lean{universal_energy_minimum}.
\begin{leanbox}[--- \lean{DipolarExciton.lean:94-96} (statement and complete proof)]
\begin{lstlisting}[style=chdeux]
theorem triangular_dipolar_optimal {d : ℝ} (hd : 0 < d) (C : Set Plane) (hC : LocallyFinite C)
    (hdens : DensityOne C) : latticeEnergy (g d) ≤ energy (g d) C :=
  (universal_energy_minimum (g d) C (dipolar_admissible hd) hC hdens).1
\end{lstlisting}
\end{leanbox}

\paragraph{Consequence 2: the inequality is not $\top\le\top$ \tierA.} An inequality in $[0,+\infty]$ whose left-hand
side equals $+\infty$ says nothing. One must therefore show that $\mathrm{latticeEnergy}(g_d)<+\infty$.
\begin{lemma}\label{ch2:lem:diffle}
For $c\ge0$ and $t>0$: $t^{-1/2}-(t+c)^{-1/2}\le\frac c2\,t^{-3/2}$. In particular $g_d(t)\le d^2t^{-3/2}$.
\end{lemma}
\begin{humanproof}
Set $a=t^{-1/2}$ and $b=(t+c)^{-1/2}$, so that $0<b\le a$, $a^2t=1$ and $b^2(t+c)=1$. Then
$a^2-b^2=a^2b^2\bigl((t+c)-t\bigr)=c\,a^2b^2$, i.e.\ $(a-b)(a+b)=c\,a^2b^2$. Since $b\le a$, $2b^2\le a^2+ab=a(a+b)$, so
$2(a-b)(a+b)\le c\,a^3(a+b)$; dividing by $a+b>0$: $a-b\le\frac c2a^3=\frac c2t^{-3/2}$.
\end{humanproof}
\noindent The sum $\sum_{a\in\Lat\setminus\{0\}}\|a\|^{-3}$ converges (Riesz exponent $3>2$, Chapter~1), hence
$\mathrm{latticeEnergy}(g_d)\le d^2\,\mathrm{latticeEnergy}(t^{-3/2})<+\infty$.
\begin{leanbox}[--- \lean{DipolarExciton.lean:119-120, 159, 167} (proofs omitted)]
\begin{lstlisting}[style=chdeux]
lemma diff_le (c : ℝ) (hc : 0 ≤ c) {t : ℝ} (ht : 0 < t) :
    t ^ (-(1 / 2 : ℝ)) - (t + c) ^ (-(1 / 2 : ℝ)) ≤ c / 2 * t ^ (-(3 / 2 : ℝ)) := by
\end{lstlisting}
\begin{lstlisting}[style=chdeux]
lemma g_le (d : ℝ) {t : ℝ} (ht : 0 < t) : g d t ≤ d ^ 2 * riesz 3 t := by
\end{lstlisting}
\begin{lstlisting}[style=chdeux]
theorem latticeEnergy_dipolar_lt_top (d : ℝ) : latticeEnergy (g d) < ⊤ := by
\end{lstlisting}
\end{leanbox}

\paragraph{Consequence 3: every density $\rho>0$ \tierA, conditional.} If $C$ has density $\rho$, then
$\sqrt\rho\,C$ has density one and $\|\sqrt\rho x-\sqrt\rho y\|^2=\rho\|x-y\|^2$: the energy of $C$ for $g$ is that of
$\sqrt\rho\,C$ for $t\mapsto g(t/\rho)$, which remains CM (each derivative gains the factor $\rho^{-r}>0$). Corollary D2
of Chapter~1 then gives the optimality of $\rho^{-1/2}\Lat$, and $g_d(t/\rho)\le d^2\rho^{3/2}t^{-3/2}$ gives finiteness.
\begin{leanbox}[--- \lean{DipolarExcitonMultipolar.lean:125-129} (statement and complete proof) and \lean{:175-176}]
\begin{lstlisting}[style=chdeux]
theorem triangular_dipolar_optimal_any_density {d : ℝ} (hd : 0 < d) {ρ : ℝ} (hρ : 0 < ρ) :
    LocallyFinite ((Real.sqrt ρ)⁻¹ • A) ∧ TriangularDensity.DensityRho ρ ((Real.sqrt ρ)⁻¹ • A) ∧
      ∀ C : Set Plane, LocallyFinite C → TriangularDensity.DensityRho ρ C →
        energy (AutoresearchDipolar.g d) ((Real.sqrt ρ)⁻¹ • A) ≤ energy (AutoresearchDipolar.g d) C :=
  TriangularDensity.triangular_optimal_any_density _ (AutoresearchDipolar.dipolar_admissible hd) hρ
\end{lstlisting}
\begin{lstlisting}[style=chdeux]
theorem energy_dipolar_any_density_lt_top {d : ℝ} (hd : 0 < d) {ρ : ℝ} (hρ : 0 < ρ) :
    energy (AutoresearchDipolar.g d) ((Real.sqrt ρ)⁻¹ • A) < ⊤ := by
\end{lstlisting}
\end{leanbox}

\paragraph{Negative controls \tierA.} (i) If the weight of the cross attraction is doubled,
$h_d(t)=t^{-1/2}-2(t+d^2)^{-1/2}$, then at $t=d^2$:
$h_d(d^2)=\frac1d-\frac{2}{\sqrt2\,d}<\frac1d-\frac{2}{2d}=0$, so $h_d$ is not admissible (\lean{h_not_admissible}).
(ii) For $c<0$ and $t=-2c>0$, we have $t+c=-c\in(0,t)$, whence $(t+c)^{-s/2}>t^{-s/2}$ and
$\mathrm{shiftDiff}_{s,c}(t)<0$ (\lean{shiftDiff_not_admissible_of_neg}). The sign of the shift matters.

\begin{warnbox}[: scope]
Finiteness is proved for $g_d$ only; for the family $\mathrm{shiftDiff}_{s,c}$ no finiteness is proved, and for $c=0$
one has $\mathrm{shiftDiff}_{s,0}\equiv0$: the conclusion is then the trivial $0\le\mathrm{energy}$. The model is
classical: rigid dipoles, no kinetic energy, no dielectric screening, no moiré potential. In the experimental density
window, real excitons form a quantum gas, not a classical crystal \tierC~\cite{ch2dip}. Paper: doi\string:10.5281/zenodo.23646103.
\end{warnbox}

% .....................................................................
\subsection{The planar Yukawa potential}\label{ch2:ssec:yuk}

\begin{physbox}[: the model \tierC]
Charged particles confined to a plane but immersed in a three-dimensional screening medium (dusty plasmas, charged
colloids) interact through the screened Coulomb potential $V(r)=e^{-\kappa r}/r$, $\kappa>0$.
In the variable $t=r^2$: $g_\kappa(t)=e^{-\kappa\sqrt t}/\sqrt t$. \textbf{This is not} the Green's function
$K_0(\kappa r)$ of the two-dimensional Debye--Hückel equation, which is not treated here~\cite{ch2yuk}.
\end{physbox}
\begin{leanbox}[--- \lean{formal_cert/Autoresearch/Yukawa.lean:28, 31}]
\begin{lstlisting}[style=chdeux]
noncomputable def phi (κ : ℝ) : ℝ → ℝ := fun t => Real.exp (-κ * Real.sqrt t)
\end{lstlisting}
\begin{lstlisting}[style=chdeux]
noncomputable def g (κ : ℝ) : ℝ → ℝ := fun t => Real.exp (-κ * Real.sqrt t) / Real.sqrt t
\end{lstlisting}
\end{leanbox}

The plan is: $g_\kappa=t^{-1/2}\cdot\varphi_\kappa$ with $\varphi_\kappa(t)=e^{-\kappa\sqrt t}$; we show that the
product of two CM functions is CM, then that $\varphi_\kappa$ is CM by strong induction. We say that $f$ has the
\emph{sign pattern up to order $n$} at $t$ if $(-1)^if^{(i)}(t)\ge0$ for $i=0,\dots,n$.

\begin{lemma}[Leibniz signs]\label{ch2:lem:leibniz}
If $f$ and $h$ are $n$ times differentiable at $t$ and both have the sign pattern up to order $n$ at $t$, then
$(-1)^n(fh)^{(n)}(t)\ge0$. Consequently, the product of two CM functions is CM.
\end{lemma}
\begin{humanproof}
Leibniz's rule: $(fh)^{(n)}=\sum_{i=0}^n\binom ni f^{(i)}h^{(n-i)}$. Write $(-1)^n=(-1)^i(-1)^{n-i}$:
\[(-1)^n(fh)^{(n)}=\sum_{i=0}^n\binom ni\,\bigl[(-1)^if^{(i)}\bigr]\,\bigl[(-1)^{n-i}h^{(n-i)}\bigr],\]
a sum of products of nonnegative numbers. Regularity and nonnegativity of a product are immediate.
\end{humanproof}
\begin{leanbox}[--- \lean{Yukawa.lean:38-41} (proof omitted; it takes about a dozen lines)]
\begin{lstlisting}[style=chdeux]
lemma mul_sign {f h : ℝ → ℝ} {t : ℝ} (r : ℕ) (hf : ContDiffAt ℝ r f t) (hh : ContDiffAt ℝ r h t)
    (Hf : ∀ i ≤ r, 0 ≤ (-1 : ℝ) ^ i * iteratedDeriv i f t)
    (Hh : ∀ i ≤ r, 0 ≤ (-1 : ℝ) ^ i * iteratedDeriv i h t) :
    0 ≤ (-1 : ℝ) ^ r * iteratedDeriv r (f * h) t := by
\end{lstlisting}
\end{leanbox}

\begin{lemma}[$t^{-s/2}$ is CM]\label{ch2:lem:riesz}
For $s>0$, $p(t)=t^{-s/2}$ satisfies $(-1)^np^{(n)}(t)=\mathrm{asc}_n(\tfrac s2)\,t^{-s/2-n}>0$.
\end{lemma}
\begin{humanproof}
This is Lemma~\ref{ch2:lem:pochder} with $c=0$ and $a=-\tfrac s2$, followed by the sign computation of Lemma~\ref{ch2:lem:sign}
with $\tfrac s2$ in place of $\tfrac12$. (Lean: \lean{TriangularRiesz.riesz_admissible}, Chapter~1.)
\end{humanproof}

\begin{proposition}[$e^{-\kappa\sqrt t}$ is CM]\label{ch2:prop:phi}
For $\kappa>0$, $\varphi_\kappa(t)=e^{-\kappa\sqrt t}$ is completely monotone on $(0,+\infty)$.
\end{proposition}
\begin{humanproof}
By the chain rule, on $(0,+\infty)$:
\[\varphi_\kappa'(t)=-\frac{\kappa}{2\sqrt t}\,e^{-\kappa\sqrt t}=-\frac\kappa2\,t^{-1/2}\,\varphi_\kappa(t).\tag{$\ast$}\]
We prove by \emph{strong induction} on $n$ the property $\mathcal P(n)$: ``for every $m\le n$ and every $t>0$,
$(-1)^m\varphi_\kappa^{(m)}(t)\ge0$''.

$\mathcal P(0)$: $\varphi_\kappa>0$.

$\mathcal P(n)\Rightarrow\mathcal P(n+1)$: only the order $m=n+1$ is new. The two sides of $(\ast)$ are equal on the
open set $(0,+\infty)$, hence so are their derivatives of all orders: $\varphi_\kappa^{(n+1)}=(\varphi_\kappa')^{(n)}
=-\frac\kappa2\,(t^{-1/2}\varphi_\kappa)^{(n)}$. By Lemma~\ref{ch2:lem:riesz}, $t^{-1/2}$ has the sign pattern at every
order; by $\mathcal P(n)$, $\varphi_\kappa$ has it up to order $n$. Lemma~\ref{ch2:lem:leibniz} gives
$(-1)^n(t^{-1/2}\varphi_\kappa)^{(n)}\ge0$, whence
\[(-1)^{n+1}\varphi_\kappa^{(n+1)}=\frac\kappa2\,(-1)^n(t^{-1/2}\varphi_\kappa)^{(n)}\ge0,\]
since $\kappa>0$. The step from $n$ to $n+1$ uses the signs of $\varphi_\kappa$ only up to order $n$: this is why strong
induction suffices and why no Faà di Bruno formula is needed. Regularity follows from the fact that
$\exp$ and $\sqrt{\,\cdot\,}$ are $C^\infty$ on $(0,+\infty)$.
\end{humanproof}
\begin{leanbox}[--- \lean{Yukawa.lean:82-83, 94-95, 121} (proofs omitted)]
\begin{lstlisting}[style=chdeux]
lemma phi_hasDerivAt (κ : ℝ) {t : ℝ} (ht : 0 < t) :
    HasDerivAt (phi κ) (-(κ / 2) * (riesz 1 t * phi κ t)) t := by
\end{lstlisting}
\begin{lstlisting}[style=chdeux]
lemma phi_sign {κ : ℝ} (hκ : 0 < κ) :
    ∀ n : ℕ, ∀ m ≤ n, ∀ t : ℝ, 0 < t → 0 ≤ (-1 : ℝ) ^ m * iteratedDeriv m (phi κ) t := by
\end{lstlisting}
\begin{lstlisting}[style=chdeux]
theorem phi_admissible {κ : ℝ} (hκ : 0 < κ) : AdmissiblePotential (phi κ) :=
\end{lstlisting}
\end{leanbox}
\noindent Note the shape of \lean{phi_sign}: the quantification ``$\forall\,m\le n$'' is exactly the strong induction of
the human proof.

\begin{theorem}[H2, admissibility of Yukawa]\label{ch2:thm:yuk}
For $\kappa>0$, $g_\kappa$ is admissible. \tierA
\end{theorem}
\begin{humanproof}
For $t>0$, $g_\kappa(t)=t^{-1/2}\varphi_\kappa(t)$: a product of two CM functions (Lemmas~\ref{ch2:lem:riesz},
\ref{ch2:lem:leibniz} and Proposition~\ref{ch2:prop:phi}). In Lean, $g_\kappa(0)=0$ by the convention $x/0=0$; but
admissibility only looks at $(0,+\infty)$, and two functions that agree on an open set have the same derivatives there
(\lean{admissible_congr}).
\end{humanproof}
\begin{leanbox}[--- \lean{Yukawa.lean:133-134} (statement and complete proof)]
\begin{lstlisting}[style=chdeux]
theorem yukawa_admissible {κ : ℝ} (hκ : 0 < κ) : AdmissiblePotential (g κ) :=
  admissible_congr (admissible_mul (riesz_admissible (s := 1) one_pos) (phi_admissible hκ)) (g_eqOn κ)
\end{lstlisting}
\end{leanbox}

\paragraph{Consequence: optimality and finite energy \tierA, conditional.} Optimality at density one
(\lean{triangular_yukawa_optimal}) is obtained as for the dipoles. For finiteness: for $x\ge0$,
$e^x\ge x^2/2$; with $x=\kappa\sqrt t$,
\[g_\kappa(t)=\frac1{e^{\kappa\sqrt t}\sqrt t}\le\frac{2}{\kappa^2\,t\sqrt t}=\frac2{\kappa^2}\,t^{-3/2},\]
and the Riesz sum with exponent $3$ converges. The hypothesis $\kappa>0$ is used here: for $\kappa=0$, $g_0=t^{-1/2}$ is the
Riesz potential with exponent $1$, whose lattice sum diverges.
\begin{leanbox}[--- \lean{Yukawa.lean:160, 175} (proofs omitted)]
\begin{lstlisting}[style=chdeux]
lemma g_le {κ : ℝ} (hκ : 0 < κ) {t : ℝ} (ht : 0 < t) : g κ t ≤ 2 / κ ^ 2 * riesz 3 t := by
\end{lstlisting}
\begin{lstlisting}[style=chdeux]
theorem latticeEnergy_yukawa_lt_top {κ : ℝ} (hκ : 0 < κ) : latticeEnergy (g κ) < ⊤ := by
\end{lstlisting}
\end{leanbox}

\paragraph{Finite mixtures \tierA.} The admissible functions form a convex cone (derivatives are linear), so every
mixture $\sum_iw_ig_{\kappa_i}$ with $w_i\ge0$, $\kappa_i>0$ is admissible, with finite energy since it is bounded by
$\bigl(\sum_iw_i\,2/\kappa_i^2\bigr)t^{-3/2}$; the result also holds at every density (\lean{YukawaMix.lean}).
\begin{leanbox}[--- \lean{formal_cert/Autoresearch/YukawaMix.lean:129-130, 138-139} (proof omitted)]
\begin{lstlisting}[style=chdeux]
noncomputable def yukawaMix {ι : Type*} (s : Finset ι) (w κ : ι → ℝ) : ℝ → ℝ :=
  fun t => ∑ i ∈ s, w i * g (κ i) t
\end{lstlisting}
\begin{lstlisting}[style=chdeux]
theorem yukawaMix_admissible {ι : Type*} (s : Finset ι) {w κ : ι → ℝ} (hw : ∀ i ∈ s, 0 ≤ w i)
    (hκ : ∀ i ∈ s, 0 < κ i) : AdmissiblePotential (yukawaMix s w κ) := by
\end{lstlisting}
\end{leanbox}

\paragraph{Negative controls \tierA.} (a) $c_\kappa(t)=e^{-\kappa\sqrt t}\cos t$ takes the value
$-e^{-\kappa\sqrt\pi}<0$ at $t=\pi$ (\lean{c_not_admissible}). (b) The signed mixture $g_1-2g_2$: at $t=\frac14$,
$\sqrt t=\frac12$ and $g_1(\tfrac14)-2g_2(\tfrac14)=2e^{-1/2}-4e^{-1}=2e^{-1}(e^{1/2}-2)<0$ since $e<4$
(\lean{signedMix_not_admissible}). These controls show that the hypotheses of the lemmas are not vacuous; they do
\emph{not} show that the triangular lattice ceases to be optimal for these functions: nothing here decides that.

\begin{warnbox}[: scope]
Classical model of point particles; no kinetic energy; not the $K_0$ kernel. The companion paper
(doi\string:10.5281/zenodo.23646105) also contains a \tierX{} numerical study with CVODE that corrects a published
truncation erratum; it supports no statement of this chapter~\cite{ch2yuk}.
\end{warnbox}

% .....................................................................
\subsection{The Rytova--Keldysh potential}\label{ch2:ssec:kel}

\begin{physbox}[: the model \tierL{} / \tierC]
In a monolayer of a two-dimensional semiconductor (MoS$_2$, WSe$_2$, \dots) suspended in vacuum, two charges
interact through the Rytova--Keldysh potential, in the form of Berkelbach, Hybertsen and Reichman~\cite{ch2bhr}:
\[V(r)=\frac{\pi}{2r_0}\,\bigl[\mathbf H_0-Y_0\bigr](r/r_0),\qquad \widehat V(k)=\frac{2\pi}{k(1+r_0k)},\]
where $\mathbf H_0$ is the Struve function, $Y_0$ the Neumann function, and $r_0>0$ a screening length
(in units $e^2/(4\pi\varepsilon_0)=1$). For $r_0\to0$ one recovers Coulomb, $2\pi/k$. The form is cited \tierL{}
(\cite{ch2bhr}, after \cite{ch2dlmf} for $\mathbf K_0=\mathbf H_0-Y_0$); the Fourier transform is \tierC.
\end{physbox}

Mathlib, at the pinned version, knows neither Struve nor Neumann functions. The idea is therefore to \emph{define} the
potential in Lean by a manifestly completely monotone integral representation --- a mixture of Gaussians in $r$, that is,
a mixture of exponentials in $t=r^2$ --- and then to prove \emph{by hand}, at Tier C, that this integral is indeed
the Struve form. The gap between the two is precisely the ``fidelity gap'' of this formalisation.

\paragraph{Step 0: the weight.} For $x\ge0$ set
\[\mathrm{erfcx}(x)=\frac{2}{\sqrt\pi}\int_0^\infty e^{-(y^2+2xy)}\,dy .\]
The change of variable $s=y+x$ gives $\mathrm{erfcx}(x)=e^{x^2}\frac2{\sqrt\pi}\int_x^\infty e^{-s^2}ds=e^{x^2}\operatorname{erfc}(x)$,
the ``scaled'' complementary error function. For $r_0>0$, the weight is $w(a)=r_0^{-1}\,\mathrm{erfcx}(\sqrt a/r_0)$.
\begin{leanbox}[--- \lean{formal_cert/Autoresearch/KeldyshCM.lean:190-191, 244, 248-249}]
\begin{lstlisting}[style=chdeux]
noncomputable def erfcx (x : ℝ) : ℝ :=
  2 / Real.sqrt Real.pi * ∫ y in Ioi (0 : ℝ), Real.exp (-(y ^ 2 + 2 * x * y))
\end{lstlisting}
\begin{lstlisting}[style=chdeux]
noncomputable def w (r0 a : ℝ) : ℝ := 1 / r0 * erfcx (Real.sqrt a / r0)
\end{lstlisting}
\begin{lstlisting}[style=chdeux]
noncomputable def keldysh (r0 t : ℝ) : ℝ :=
  ∫ a in Ioi (0 : ℝ), w r0 a / (2 * a) * Real.exp (-t / (4 * a))
\end{lstlisting}
\end{leanbox}
\noindent In other words, the Lean object is
\begin{equation}\label{ch2:eq:kmix}
\mathcal K_{r_0}(t)=\int_0^\infty\frac{w(a)}{2a}\,e^{-t/(4a)}\,da .
\end{equation}

\begin{lemma}[bounds on the weight]\label{ch2:lem:w}
For $x\ge0$, $0<\mathrm{erfcx}(x)\le1$; for $x>0$, $\mathrm{erfcx}(x)\le\frac1{\sqrt\pi\,x}$. Hence, for $a>0$,
$0<w(a)\le\frac1{r_0}$ and $w(a)\le\frac1{\sqrt\pi\sqrt a}$.
\end{lemma}
\begin{humanproof}
The integrand is continuous and strictly positive, so $\mathrm{erfcx}(x)>0$. For $y>0$ and $x\ge0$,
$y^2+2xy\ge y^2$, so the integral is at most $\int_0^\infty e^{-y^2}dy=\frac{\sqrt\pi}2$, whence $\mathrm{erfcx}\le1$. If
$x>0$, bound from below by $2xy$ instead: $\int_0^\infty e^{-2xy}dy=\frac1{2x}$, whence
$\mathrm{erfcx}(x)\le\frac2{\sqrt\pi}\cdot\frac1{2x}$. For $w$, substitute $x=\sqrt a/r_0$.
\end{humanproof}

\paragraph{Step 1: a generic Laplace-mixture lemma.} Let $(\alpha,\mu)$ be a measure space, $k,m$ measurable,
$k\ge0$ and $m\ge0$ almost everywhere, and assume $\int k^ne^{-ck}m\,d\mu<\infty$ for every $n$ and every $c>0$. Set
$G_n(t)=\int(-k)^ne^{-tk}m\,d\mu$.
\begin{theorem}[Laplace mixture]\label{ch2:thm:laplace}
$g=G_0$ is admissible, and $g^{(n)}(t)=G_n(t)$ for $t>0$. \tierA
\end{theorem}
\begin{humanproof}
Fix $t>0$ and let $s$ range over $(t/2,+\infty)$, a neighbourhood of $t$. The derivative in $s$ of $(-k)^ne^{-sk}m$ is
$(-k)^{n+1}e^{-sk}m$, whose absolute value satisfies $k^{n+1}e^{-sk}m\le k^{n+1}e^{-(t/2)k}m$ since $s\ge t/2$ and
$k\ge0$. The bound does not depend on $s$ and is integrable (hypothesis with $c=t/2$). The theorem on differentiation
under the integral sign (dominated convergence) gives $G_n'(t)=G_{n+1}(t)$. By induction, $G_0$ is of class $C^N$ for
every $N$ and $g^{(n)}=G_n$. Finally $(-1)^nG_n(t)=\int k^ne^{-tk}m\,d\mu\ge0$, since the integrand is nonnegative almost
everywhere.
\end{humanproof}
\begin{uniremark}\tierL{} The step used above is the standard theorem on differentiation under the integral sign: let
$U\subset\R$ be open, $(X,\mu)$ a measure space, $F\colon U\times X\to\R$ with $F(s,\cdot)$ integrable for every $s\in U$, $s\mapsto F(s,x)$
differentiable for almost every $x$, and $|\partial_sF(s,x)|\le H(x)$ for all $s\in U$ and almost every $x$, with $H$
integrable and independent of $s$. Then $s\mapsto\int F(s,x)\,d\mu(x)$ is differentiable on $U$ with derivative
$\int\partial_sF(s,x)\,d\mu(x)$. It follows from the mean value theorem and Lebesgue's dominated convergence theorem
(standard).
\end{uniremark}
\begin{leanbox}[--- \lean{KeldyshCM.lean:31-36, 134-135} (proof omitted)]
\begin{lstlisting}[style=chdeux]
structure MixtureData (μ : Measure α) (k m : α → ℝ) : Prop where
  k_meas : Measurable k
  m_meas : Measurable m
  k_nonneg : ∀ᵐ x ∂μ, 0 ≤ k x
  m_nonneg : ∀ᵐ x ∂μ, 0 ≤ m x
  integrable : ∀ (n : ℕ) (c : ℝ), 0 < c → Integrable (fun x => k x ^ n * Real.exp (-(c * k x)) * m x) μ
\end{lstlisting}
\begin{lstlisting}[style=chdeux]
theorem laplace_admissible (h : MixtureData μ k m) :
    AdmissiblePotential (fun t => ∫ x, Real.exp (-(t * k x)) * m x ∂μ) := by
\end{lstlisting}
\end{leanbox}

This theorem is re-exported under the name \lean{ALEXIS} (Chapter~5).

\paragraph{Step 2: the Keldysh instance.} Take Lebesgue measure on $(0,+\infty)$, $k(a)=\frac1{4a}$ and
$m(a)=\frac{w(a)}{2a}$. The moments must be checked: $F_n(a)=(4a)^{-n}e^{-c/(4a)}\frac{w(a)}{2a}$ is integrable.
Split at $a=1$. On $(0,1]$, with $u=\frac1{4a}$: $F_n=2w(a)\,u^{n+1}e^{-cu}$, and $e^{cu}\ge(cu)^N/N!$ gives
$u^Ne^{-cu}\le N!/c^N$; $F_n$ is bounded by $2r_0^{-1}(n+1)!/c^{n+1}$ on an interval of finite length. On
$(1,+\infty)$, $u\le\frac14$, so $F_n\le\frac{w(a)}{2a}\le\frac1{2\sqrt\pi}a^{-3/2}$ (Lemma~\ref{ch2:lem:w}), and
$\int_1^\infty a^{-3/2}da=2$. Theorem~\ref{ch2:thm:laplace} applies.
\begin{leanbox}[--- \lean{KeldyshCM.lean:354} (proof omitted)]
\begin{lstlisting}[style=chdeux]
theorem keldysh_admissible {r0 : ℝ} (hr0 : 0 < r0) : AdmissiblePotential (keldysh r0) := by
\end{lstlisting}
\end{leanbox}

\paragraph{Step 3: the statement is not that of the zero function \tierA.} In Lean, the Bochner integral of a
non-integrable function is $0$, and $0$ is admissible: admissibility alone could be vacuous. But for $t>0$
the integrand of~\eqref{ch2:eq:kmix} is strictly positive and integrable, so $\mathcal K_{r_0}(t)>0$.
\begin{leanbox}[--- \lean{KeldyshCM.lean:379} (proof omitted)]
\begin{lstlisting}[style=chdeux]
theorem keldysh_pos {r0 t : ℝ} (hr0 : 0 < r0) (ht : 0 < t) : 0 < keldysh r0 t := by
\end{lstlisting}
\end{leanbox}

\paragraph{Step 4: corollaries, and a caveat \tierA{} / \tierC.} The product with $e^{-\varepsilon t}$
($\varepsilon\ge0$) remains admissible (Lemma~\ref{ch2:lem:leibniz}), and the upstream theorem gives the two inequalities
\lean{triangular_keldysh_optimal} and \lean{triangular_keldyshReg_optimal}.
\begin{leanbox}[--- \lean{KeldyshCM.lean:390-392} (statement and complete proof)]
\begin{lstlisting}[style=chdeux]
theorem triangular_keldysh_optimal {r0 : ℝ} (hr0 : 0 < r0) (C : Set Plane) (hC : LocallyFinite C)
    (hdens : DensityOne C) : latticeEnergy (keldysh r0) ≤ energy (keldysh r0) C :=
  (universal_energy_minimum (keldysh r0) C (keldysh_admissible hr0) hC hdens).1
\end{lstlisting}
\end{leanbox}
\begin{warnbox}[: $\top\le\top$]
$V(r)\sim1/r$ at large distance, and $\sum_{0<|x|\le R}1/|x|$ grows like $R$ on a planar lattice. One therefore
\emph{expects} $\mathrm{latticeEnergy}(\mathcal K_{r_0})=+\infty$: the first inequality would then read $\top\le\top$
\tierC. This is neither proved nor refuted in Lean. For $\varepsilon>0$, the regularised sum should be finite \tierC,
but \textbf{this finiteness is not proved}. The physically relevant energy of a Wigner crystal, neutralised
by a uniform background (``jellium''), is covered by neither of the two statements~\cite{ch2kel}.
\end{warnbox}

\paragraph{Step 5: the link with the Struve form \tierC.}
\begin{theorem}[not machine-checked]\label{ch2:thm:struve}
For $r_0>0$ and $r>0$: $V(r)=\mathcal K_{r_0}(r^2)$.
\end{theorem}
\begin{humanproof}
All integrands are nonnegative: every exchange of integrals is justified by Tonelli's theorem.

\emph{(a) A Laplace form in $r$.} Formula DLMF 11.5.2 with $\nu=0$ gives
$\mathbf K_0(z)=\frac2\pi\int_0^\infty e^{-zu}(1+u^2)^{-1/2}du$ for $\operatorname{Re}z>0$ \tierL~\cite{ch2dlmf}. With
$z=r/r_0$: $V(r)=\frac1{r_0}\int_0^\infty e^{-ru/r_0}(1+u^2)^{-1/2}du$. This already suffices to see that $V$ is CM \emph{in $r$}.

\emph{(b) Subordination.} For $\rho>0$ and $t>0$,
$e^{-\rho\sqrt t}=\int_0^\infty e^{-ta}\,\frac{\rho}{2\sqrt\pi}\,a^{-3/2}e^{-\rho^2/(4a)}\,da$ (Laplace transform of
the Lévy density, a classical formula used without citation in the source paper). Take $\rho=u/r_0$, $t=r^2$,
insert into (a) and exchange the integrals:
\[V(r)=\int_0^\infty e^{-ta}\,\frac{a^{-3/2}}{2\sqrt\pi\,r_0^2}\left[\int_0^\infty\frac{u}{\sqrt{1+u^2}}e^{-\lambda u^2}du\right]da,
\qquad\lambda=\frac1{4r_0^2a}.\]

\emph{(c) The inner integral is elementary.} With $v=\sqrt{1+u^2}$, $dv=u(1+u^2)^{-1/2}du$:
\[\int_0^\infty\frac{u}{\sqrt{1+u^2}}e^{-\lambda u^2}du=e^{\lambda}\int_1^\infty e^{-\lambda v^2}dv
=\frac12\sqrt{\frac\pi\lambda}\,\mathrm{erfcx}(\sqrt\lambda).\]
With $\sqrt{\pi/\lambda}=2r_0\sqrt{\pi a}$, one obtains $V(r)=\int_0^\infty e^{-ta}\rho(a)\,da$ where
$\rho(a)=\dfrac{\mathrm{erfcx}\bigl(1/(2r_0\sqrt a)\bigr)}{2r_0\,a}$.

\emph{(d) Change of variable.} Set $a=\frac1{4b}$, whence $|da|=\frac{db}{4b^2}$ and $e^{-ta}=e^{-t/(4b)}$. Then
$\rho(a)\,da=\frac{4b\,\mathrm{erfcx}(\sqrt b/r_0)}{2r_0}\cdot\frac{db}{4b^2}=\frac{w(b)}{2b}\,db$, so
$V(r)=\int_0^\infty\frac{w(b)}{2b}e^{-r^2/(4b)}db=\mathcal K_{r_0}(r^2)$.
\end{humanproof}
\noindent Numerical checks at 30 digits (mpmath) confirm each step, with two planted errors
detected \tierX~\cite{ch2kel}; they do not replace the proof. The reader must accept
Theorem~\ref{ch2:thm:struve} to transfer the Lean statement, which concerns~\eqref{ch2:eq:kmix}, to the Struve form.
Paper: doi\string:10.5281/zenodo.23646107.

% ---------------------------------------------------------------------
\section{The generalisation: a CM/Bernstein library}\label{ch2:sec:lib}

The three preceding proofs share a family resemblance: a difference of translates (dipoles), a product and a
composition with $\sqrt t$ (Yukawa), a Laplace mixture (Keldysh). The P3 library~\cite{p3} (doi\string:10.5281/zenodo.23653842)
isolates these mechanisms once and for all. Gate summary: \textbf{130 theorems} declared, each with its
\texttt{\#print axioms} line, all within the allowed list, no warnings; an audit by a separate AI agent
rebuilt the 11 files and 5 prerequisite modules from source and found 0 off-list constants among 167
\tierA. No human audit; the 20 \texttt{LEDGER\_UNAUDITED\_TIER\_A} flags remain open.

\subsection{Two definitions and a bridge}
\begin{definition}[Bernstein function]
$\psi\colon(0,+\infty)\to\R$ is a \emph{Bernstein function} if it is $C^\infty$, nonnegative ($\psi\ge0$), and $\psi'$ is CM.
Examples: $at+b$ ($a,b\ge0$), $t^\alpha$ for $0\le\alpha\le1$, in particular $\sqrt t$; $1-e^{-at}$; $\ln(1+t)$.
\end{definition}
\begin{uniremark}\tierL{} Bernstein functions admit a Lévy--Khintchine representation: $\psi$ is a Bernstein function if
and only if $\psi(t)=a+bt+\int_{(0,\infty)}(1-e^{-tu})\,\nu(du)$ for constants $a,b\ge0$ and a positive measure $\nu$ on
$(0,\infty)$ with $\int\min(1,u)\,\nu(du)<\infty$~\cite{ch2ssv}. For example, $\sqrt t=\frac1{2\sqrt\pi}\int_0^\infty(1-e^{-tu})\,u^{-3/2}\,du$.
This representation is not formalised here; the library works only with the derivative-sign definition above.
\end{uniremark}
\begin{leanbox}[--- \lean{formal_cert/Programmes/CompletelyMonotone/Basic.lean:28-30, 33-34} and \lean{Bernstein.lean:24-25}]
\begin{lstlisting}[style=chdeux]
def IsCompletelyMonotoneOn (f : ℝ → ℝ) : Prop :=
  ContDiffOn ℝ (⊤ : ℕ∞) f (Set.Ioi 0) ∧
  ∀ (n : ℕ) (t : ℝ), 0 < t → 0 ≤ (-1 : ℝ) ^ n * iteratedDeriv n f t
\end{lstlisting}
\begin{lstlisting}[style=chdeux]
theorem isCompletelyMonotoneOn_iff_admissible {f : ℝ → ℝ} :
    IsCompletelyMonotoneOn f ↔ OAI.AtomicTriangular.AdmissiblePotential f := by
\end{lstlisting}
\begin{lstlisting}[style=chdeux]
def IsBernsteinOn (g : ℝ → ℝ) : Prop :=
  ContDiffOn ℝ (⊤ : ℕ∞) g (Set.Ioi 0) ∧ (∀ t : ℝ, 0 < t → 0 ≤ g t) ∧ IsCompletelyMonotoneOn (deriv g)
\end{lstlisting}
\end{leanbox}
\noindent The bridge \lean{isCompletelyMonotoneOn_iff_admissible} is an \emph{equivalence} \tierA: the clause $g\ge0$ of
the upstream is the case $n=0$ of the third clause. Every CM function of the library can therefore be used directly in
\lean{universal_energy_minimum}.

\subsection{The closure lemmas}
Write $h_n(t)=(-1)^nf^{(n)}(t)$ for $f$ CM.
\begin{enumerate}[label=(\roman*),itemsep=2pt]
\item \emph{Cone.} Derivatives are linear: $cf$ ($c\ge0$), $f+g$, and every finite nonnegative combination of
CM functions is CM. Nonnegative constants are CM.
\item \emph{Product.} By the Leibniz lemma~\ref{ch2:lem:leibniz}, $fg$ is CM.
\item \emph{Derivative.} $h_n'=(-1)^nf^{(n+1)}=-h_{n+1}\le0$, so each $h_n$ is nonincreasing (mean value theorem). In
particular $-f'$ is CM, since $(-1)^n(-f')^{(n)}=h_{n+1}\ge0$.
\item \emph{Translate.} For $c\ge0$, $t\mapsto f(t+c)$ is CM: its derivatives are $f^{(n)}(t+c)$ and $t+c>0$.
\item \emph{Difference of translates.} For $c\ge0$, $f(t)-f(t+c)$ is CM, since
$(-1)^n\bigl[f(t)-f(t+c)\bigr]^{(n)}=h_n(t)-h_n(t+c)\ge0$ because $h_n$ is nonincreasing. No integral representation is used.
\end{enumerate}
\begin{leanbox}[--- \lean{Basic.lean:160-161} (proof omitted; it takes ten lines)]
\begin{lstlisting}[style=chdeux]
theorem sub_comp_add_const (hf : IsCompletelyMonotoneOn f) {c : ℝ} (hc : 0 ≤ c) :
    IsCompletelyMonotoneOn (fun t => f t - f (t + c)) := by
\end{lstlisting}
\end{leanbox}
\noindent Item (v) is the abstract version of the dipolar proof: the Pochhammer computation is replaced there by the
mere monotonicity of $h_n$.

\subsection{The subordination theorem}
\begin{theorem}[subordination]\label{ch2:thm:subord}
If $f$ is CM and $\psi$ is a Bernstein function that is strictly positive on $(0,+\infty)$, then $f\circ\psi$ is
CM. \tierA{}; proof below \tierC.
\end{theorem}
\begin{humanproof}
Regularity comes from the stability of $C^\infty$ under composition ($\psi$ maps $(0,+\infty)$ into $(0,+\infty)$).
For the signs, we prove by induction on $n$ the property
\[\mathcal Q(n)\ :\ \text{``for \emph{every} CM function } f\text{ and every } m\le n,\ (-1)^m(f\circ\psi)^{(m)}\ge0\text{ on }(0,+\infty)\text{''.}\]
$\mathcal Q(0)$: $f(\psi(t))\ge0$ since $\psi(t)>0$ and $f\ge0$.

$\mathcal Q(n)\Rightarrow\mathcal Q(n+1)$: let $f$ be CM; the orders $m\le n$ are covered by $\mathcal Q(n)$. For
$m=n+1$, set $F=-f'$, which is CM by (iii). The chain rule gives
$(f\circ\psi)'=(f'\circ\psi)\,\psi'=-(F\circ\psi)\,\psi'$, so
\[(-1)^{n+1}(f\circ\psi)^{(n+1)}=(-1)^n\bigl[(F\circ\psi)\,\psi'\bigr]^{(n)}\ge0\]
by the Leibniz lemma, applied to $u=F\circ\psi$ (whose signs up to order $n$ are given by $\mathcal Q(n)$
\emph{applied to $F$}) and to $v=\psi'$ (CM by the definition of a Bernstein function). It is because $\mathcal Q(n)$
quantifies over \emph{all} CM functions that the induction closes without Faà di Bruno's formula.
\end{humanproof}
\begin{leanbox}[--- \lean{Bernstein.lean:49-51} (proof omitted) and \lean{:85-87} (statement and complete proof)]
\begin{lstlisting}[style=chdeux]
theorem comp_sign (hg : IsBernsteinOn g) (hpos : ∀ t, 0 < t → 0 < g t) :
    ∀ n : ℕ, ∀ f : ℝ → ℝ, IsCompletelyMonotoneOn f → ∀ m ≤ n, ∀ t : ℝ, 0 < t →
      0 ≤ (-1 : ℝ) ^ m * iteratedDeriv m (fun t => f (g t)) t := by
\end{lstlisting}
\begin{lstlisting}[style=chdeux]
theorem IsCompletelyMonotoneOn.comp_bernstein {f g : ℝ → ℝ} (hf : IsCompletelyMonotoneOn f)
    (hg : IsBernsteinOn g) (hpos : ∀ t, 0 < t → 0 < g t) : IsCompletelyMonotoneOn (fun t => f (g t)) :=
  ⟨hf.1.comp hg.1 (fun s hs => hpos s hs), fun n t ht => hg.comp_sign hpos n f hf n le_rfl t ht⟩
\end{lstlisting}
\end{leanbox}
\noindent The Lean statement \lean{comp_sign} carries ``$\forall f$'' \emph{inside} the induction on $n$: it is
literally the property $\mathcal Q(n)$.

\begin{remark}[the hypothesis $\psi>0$ is necessary \tierA]
$f$ is constrained only on $(0,+\infty)$, and $\psi\equiv0$ is a Bernstein function; then $f\circ\psi$ is the ``junk''
constant $f(0)$, which may be negative. The control \lean{junk_comp_zero_not_cm} realises this case. The theorem is
classical in spirit (``Bochner subordination''); the source paper did not locate its number in the
reference monograph~\cite{ch2ssv}. The contribution is the machine-checked proof on $(0,+\infty)$.
\end{remark}

\paragraph{Instances.} $\psi(t)=t^\alpha$ ($0\le\alpha\le1$) is a Bernstein function: $\psi'=\alpha t^{-(1-\alpha)}$ is CM
by Lemma~\ref{ch2:lem:riesz}. Hence, with $f(u)=e^{-au}$: $e^{-at^\alpha}$ is CM (stretched exponential), and with
$\alpha=\frac12$: $e^{-\kappa\sqrt t}$ is CM for every $\kappa\ge0$, in one line.
\begin{leanbox}[--- \lean{Bernstein.lean:160-161} (statement and complete proof)]
\begin{lstlisting}[style=chdeux]
theorem phi_cm {κ : ℝ} (hκ : 0 ≤ κ) : IsCompletelyMonotoneOn (AutoresearchYukawa.phi κ) :=
  ((exp_neg_mul hκ).comp_sqrt).congr' fun t _ => by simp only [AutoresearchYukawa.phi, neg_mul]
\end{lstlisting}
\end{leanbox}
\noindent Compare with the strong induction of Proposition~\ref{ch2:prop:phi}: that one required $\kappa>0$ and an
ad hoc induction; the library covers $\kappa\ge0$.

\paragraph{Laplace mixtures, weakened hypothesis.} The theorem \lean{laplace} of \lean{Laplace.lean:49} no longer
assumes that all moments are finite, but only that $e^{-ck}m$ is integrable for every $c>0$: the inequality
$u^n\le n!\,(c/2)^{-n}e^{(c/2)u}$ ($u\ge0$) gives $k^ne^{-ck}m\le\frac{n!}{(c/2)^n}e^{-(c/2)k}m$, whence all the
moments. This is the \emph{easy} direction of the Bernstein--Widder theorem (a Laplace transform of a positive
measure is CM); the hard direction (every CM function is such a transform~\cite{widder}) is
\emph{not} proved.
\begin{uniremark}\tierL{} The full statement (Bernstein's theorem, also called the Bernstein--Widder theorem): a
function $f\colon(0,\infty)\to\R$ is completely monotone if and only if $f(t)=\int_{[0,\infty)}e^{-tu}\,\mu(du)$ for
a positive Borel measure $\mu$ on $[0,\infty)$ such that the integral is finite for every $t>0$; the measure $\mu$ is
then unique~\cite{bernstein,widder}. The measure need not be finite: $1/t=\int_0^\infty e^{-tu}\,du$. The preceding paragraph concerns only the ``if''
direction; the ``only if'' direction is not proved in this book.
\end{uniremark}

\subsection{The three potentials become corollaries}
\begin{leanbox}[--- \lean{Corollaries.lean:24-25, 28-30, 38}]
\begin{lstlisting}[style=chdeux]
theorem yukawa_cm {κ : ℝ} (hκ : 0 ≤ κ) : IsCompletelyMonotoneOn (AutoresearchYukawa.g κ) :=
  ((riesz_cm zero_le_one).mul (phi_cm hκ)).congr (AutoresearchYukawa.g_eqOn κ)
\end{lstlisting}
\begin{lstlisting}[style=chdeux]
theorem dipolar_cm (d : ℝ) : IsCompletelyMonotoneOn (AutoresearchDipolar.g d) :=
  (((rpow_neg (by norm_num : (0 : ℝ) ≤ 1 / 2)).sub_comp_add_const (sq_nonneg d)).const_mul
    zero_le_two).congr' fun _ _ => rfl
\end{lstlisting}
\begin{lstlisting}[style=chdeux]
theorem keldysh_cm {r0 : ℝ} (hr0 : 0 < r0) : IsCompletelyMonotoneOn (AutoresearchKeldysh.keldysh r0) := by
\end{lstlisting}
\end{leanbox}
\noindent Reading: Yukawa $=$ Riesz $\times$ ($\exp\circ\sqrt{\ }$); dipole $=2\times$ (difference of translates of
$t^{-1/2}$); Keldysh $=$ Laplace mixture (eight-line proof, omitted).

\begin{warnbox}[: what the library does not shorten]
The phrase ``one-line corollaries'' is only half true. The P3 paper \emph{measured} the dependency closure of the three
results H1, H2, H3 taken together: it \emph{increases} when one goes through the library. Only the \emph{marginal} cost
of a new potential becomes small. For Keldysh, the short proof still relies on the heavy bounds of
\lean{keldysh_data} (erfcx, integrability) from the original module: the weakened hypothesis is
\emph{available} but not \emph{exercised}. Mathlib already contains a neighbouring notion, \lean{AbsolutelyMonotoneOn}
(without alternating signs); the library is not ready for Mathlib~\cite{p3}.
\end{warnbox}

\paragraph{New consequences \tierA, conditional for optimality.} Soft-core Yukawa
$g_\kappa(t+b^2)$ (translate), stretched exponential $e^{-at^\alpha}$, screened dipolar exciton
$2(g_\kappa(t)-g_\kappa(t+d^2))$, continuous mixtures of Yukawa potentials. An instructive example, the quadrupoles $+1,-2,+1$ in
the layers $0,d,2d$ ($c=d^2$):
\begin{itemize}
\item the \emph{field} of a quadrupole on a test charge, $f(t)-2f(t+c)+f(t+4c)$ with $f=t^{-1/2}$, is \emph{not}
CM: at $t=3$, $c=1$, it equals $\frac1{\sqrt3}-1+\frac1{\sqrt7}<0$ (\lean{quadrupole_not_cm});
\item the \emph{interaction} of two aligned quadrupoles counts the pairs ($1+4+1=6$ at $\Delta z=0$, $4\times(-2)=-8$ at
$|\Delta z|=d$, $2\times1=2$ at $|\Delta z|=2d$) and equals $2\bigl(3f(t)-4f(t+c)+f(t+4c)\bigr)$ \tierC{} for the counting;
it \emph{is} CM for every CM $f$ (Exercise~\ref{ch2:ex:quad}).
\end{itemize}
\begin{leanbox}[--- \lean{Applications.lean:172, 175-176} (proof omitted)]
\begin{lstlisting}[style=chdeux]
noncomputable def quadrupolarPair (f : ℝ → ℝ) (c t : ℝ) : ℝ := 2 * (3 * f t - 4 * f (t + c) + f (t + 4 * c))
\end{lstlisting}
\begin{lstlisting}[style=chdeux]
theorem quadrupolarPair_cm {f : ℝ → ℝ} (hf : IsCompletelyMonotoneOn f) {c : ℝ} (hc : 0 ≤ c) :
    IsCompletelyMonotoneOn (quadrupolarPair f c) := by
\end{lstlisting}
\end{leanbox}
\noindent For these new potentials, no finiteness of the optimal energy is proved: for slowly decaying tails,
the optimality inequality may read $\top\le\top$.

% ---------------------------------------------------------------------
\section{A neighbouring problem: the Sandier--Serfaty conjecture}\label{ch2:sec:ss}

\begin{warnbox}[: status]
\textbf{The Sandier--Serfaty conjecture is not proved in this book.} A manuscript in the openai/math release
claims a proof of it; this manuscript is AI-generated and not peer reviewed. What we have
done --- reproduce a certificate, audit a reduction, write a Lean statement, formalise a small part --- constitutes
neither a proof, nor a verification, nor a peer review.
\end{warnbox}

\subsection{What the conjecture says}\label{ch2:ssec:ssconj}
\begin{physbox}[: from vortices to lattices \tierC]
In a type-II superconductor in a magnetic field, the field penetrates in the form of quantised flux tubes,
the \emph{vortices}, which arrange themselves in a lattice (the Abrikosov lattice, observed to be triangular). Sandier and Serfaty
derived from the Ginzburg--Landau model a \emph{renormalised energy} $W$ that governs the position of the vortices at
next order~\cite{ch2ss} \tierL. The link with vortex physics is a \tierC{} gloss here.
\end{physbox}
Here is the formulation of the manuscript, which follows Sandier--Serfaty \tierL. Let $\Lambda\subset\R^2$ be locally
finite, of bounded density, and $\nu_\Lambda=\sum_{p\in\Lambda}\delta_p$. A locally integrable field $E:\R^2\to\R^2$ is
\emph{admissible} ($E\in\mathcal A_1$) if, in the sense of distributions,
\[\operatorname{div}E=2\pi(\nu_\Lambda-1),\qquad\operatorname{curl}E=0 .\]
Each point of $\Lambda$ is a unit charge; the ``$-1$'' is a uniform neutralising background. Near a charge,
$|E|^2\sim1/|x-p|^2$ is not integrable: one removes disks of radius $\eta$ and compensates,
\[W(E,\chi_R)=\lim_{\eta\to0}\Bigl[\tfrac12\int_{\R^2\setminus\bigcup_pB(p,\eta)}\chi_R|E|^2+\pi\ln\eta\sum_{p\in\Lambda}\chi_R(p)\Bigr],
\qquad W(E)=\limsup_{R\to\infty}\frac{W(E,\chi_R)}{4R^2},\]
where $\chi_R$ is a smooth cutoff supported in the square $[-R,R]^2$. The sign $+\pi\ln\eta$ is the right one:
$\frac12\int_{\eta<|x-p|<\rho}|x-p|^{-2}dx=\pi\ln\rho-\pi\ln\eta$.
\begin{itemize}
\item \textbf{Conjecture} (Sandier--Serfaty 2012, Conjecture~1): $W(E)\ge W(E_\triangle)$ for every $E\in\mathcal A_1$,
where $E_\triangle$ is the periodic field of the triangular lattice $\Lambda_\triangle$ of covolume one \tierL.
\item Sandier--Serfaty prove it \emph{among lattices} (their Theorem~2) and write that minimisation in
general remains open \tierL~\cite{ch2ss}.
\item Bétermin--Sandier show that it is equivalent to a value conjectured by Brauchart--Hardin--Saff for the
next-order term of the optimal logarithmic energy of $N$ points on the sphere~\cite{ch2bs,ch2bhs} \tierL, and give
$W(E_\triangle)=\pi\ln\frac{2\sqrt2\,\pi}{\sqrt3\,\Gamma(1/3)^3}\approx-4.1504128$ \tierL.
\end{itemize}

\paragraph{What the openai/math manuscript claims.} A proof~\cite{ch2oaiss}, via (1) a reduction from the plane to
square tori $\R^2/(\sqrt n\,\Z^2)$, (2) a trial potential glued on the Voronoi sectors, \emph{exactly
calibrated} to the triangular lattice, and (3) a finite certificate in exact interval arithmetic (144-bit integers),
organised in twelve stages. This manuscript is AI-generated and not peer reviewed.

\subsection{P1: reproduce and audit, without proving}\label{ch2:ssec:p1}
\begin{itemize}
\item \textbf{Reproduction of the certificate \tierX.} The unchanged C++ program was recompiled and run on a different
toolchain and a different machine (g++~11.4); running the certificate took about 478~s; the final line is
\texttt{OK} and the 22 printed integer pairs are identical, 22 out of 22, to the original record.
\emph{Negative controls}: three one-line weakened versions; two are rejected by the program's assertions,
\textbf{the third one passes} (the outcomes of the last two had been predicted in writing before the computations
finished). The reproduction shows that the program behaves as recorded; it tests none of the analytic bounds it is
supposed to certify.
\item \textbf{Audit of the reduction \tierC.} An AI agent checked Sections 2 to 5 of the manuscript step by step
(31 numbered steps and 14 lines of the certificate step): \emph{no step was classified as a gap}, with declared
limits (the spline pieces outside $[p,R_0]$, the values of the grid step and trust in the
execution are not independently checked). This is neither a proof nor a peer review.
\item \textbf{Locked Lean statement.} The conjecture is written as a Lean proposition (\lean{Prop}),
\textbf{not proved}, whose statement region is locked by fingerprint (\texttt{statement\_lock.py}).
\end{itemize}
\begin{leanbox}[--- \lean{formal_cert/Programmes/SandierSerfatyStatement.lean:153-160} (a definition, not a theorem)]
\begin{lstlisting}[style=chdeux]
def SandierSerfatyStatement : Prop :=
  (∃ χ, IsSquareCutoffFamily χ) ∧
    ∃ E₀ : Plane → Plane, Admissible triLattice E₀ ∧ IsPeriodic triLattice E₀ ∧
      ∀ χ : ℝ → Plane → ℝ, IsSquareCutoffFamily χ →
        ∃ w₀ : ℝ → ℝ, (∀ R : ℝ, 1 < R → HasLocalEnergy triLattice E₀ (χ R) (w₀ R)) ∧
          WSq w₀ ≠ ⊤ ∧ WSq w₀ ≠ ⊥ ∧
          ∀ (Λ : Set Plane) (E : Plane → Plane) (w : ℝ → ℝ), Admissible Λ E →
            (∀ R : ℝ, 1 < R → HasLocalEnergy Λ E (χ R) (w R)) → WSq w₀ ≤ WSq w
\end{lstlisting}
\end{leanbox}
\noindent Only ten non-degeneracy lemmas are \tierAs{}: the admissible class is nonempty (the field $E=-\pi x$,
with no charge), the background constant $2\pi$ matters (with $0$ or $-2\pi$ in its place, this field is no longer
admissible), the cutoff classes are nonempty. Warning: the witness of \lean{competitor_nonvacuous} is this
field $-\pi x$, whose energy is $+\infty$ \tierC{}; it shows that the hypotheses on a competitor are
\emph{satisfiable}, not that the inequality constrains a competitor. The existence of the limit in $\eta$ is \emph{assumed}
(\lean{HasLocalEnergy}) and the triangular minimiser is an existential witness: these are declared deviations from the
manuscript's formulation~\cite{ch2ssa}. P1 paper: doi\string:10.5281/zenodo.23653827.

\subsection{A1: formalise a part, and measure the rest}\label{ch2:ssec:a1}
\paragraph{Interval arithmetic in Lean.} An interval is a pair of integers $(u,v)$ read as $[u/Q,v/Q]$ with
$Q=2^{144}$. Products are rounded \emph{outwards}: lower bound by $\lfloor\cdot/Q\rfloor$, upper
bound by $\lceil\cdot/Q\rceil$, so that the true real result is never lost. Toy example with $Q=100$:
$[1.20, 1.30]\times[0.50, 0.60]$, i.e.\ $(120,130)\times(50,60)$, has endpoint products
$6000,7200,6500,7800$, whence $(\lfloor6000/100\rfloor,\lceil7800/100\rceil)=(60,78)=[0.60, 0.78]$~\cite{ch2ssl}.

\paragraph{Two verification levels} (an axis distinct from Tiers A/L/C/X):
\textbf{L1} (kernel re-execution): the Lean transcription of a stage, evaluated by the kernel with
\lean{decide +kernel} (and not \lean{native_decide}, which adds an axiom and trusts compiled code), reproduces
bit for bit the integers printed by the C++. \textbf{L2} (real statement): a Lean theorem asserts that a \emph{real
number} belongs to the interval.

\begin{center}\small
\begin{tabular}{@{}p{0.27\linewidth}p{0.4\linewidth}p{0.29\linewidth}@{}}\toprule
Result (saved logs) & Value & Tier\\\midrule
Stages complete at L1 & 4 of 12 (\lean{init_constants}, \lean{bcons/eps}, \lean{calcLam}, \lean{calcK}) & \tierAs{} for \lean{init_constants}, \lean{bcons/eps}; \tierC{} for \lean{calcLam}, \lean{calcK} (no independent recompilation)\\
\quad of which conditional & \lean{calcLam}, \lean{calcK} (tables as premises), \lean{bcons} (coefficients $c_m$) & \\
Stages complete at L2 & 1 of 12 (\lean{init_constants}: $\ln2,\pi,p,R_0,a$ real) & \tierAs\\
Share of interval operations re-executed & $1\,685\,129$ of $1\,238\,829\,689$, i.e.\ $0.136\,\%$ & \tierX\\
Planted false statements rejected & 11 & \tierX\\
Projected cost of the rest at L1 & 854 to 2\,240 core-hours & \tierX\\\bottomrule
\end{tabular}
\end{center}
The 22 printed pairs and the final assertions of the C++ are in addition checked as arithmetic on the exact
endpoints (level ``E'', weaker), except the bound test of the \lean{core_test} stage; the table-separation test
is checked only window by window. Independent verification: the audit by a second AI agent rebuilt
13 Lean modules from source, but \emph{not} \lean{CalcK.lean}, \lean{CalcLam.lean} nor 28 chunk files;
the L1 results of \lean{calcLam} and \lean{calcK} therefore have not had a producer${}\ne{}$verifier pass
\tierC~\cite{ch2ssl}. The heaviest remaining work
is level L2 of the analytic content (radial functions, lattice coefficients, quadrature remainders), then
30 further steps of the reduction: a project of several months \tierC~\cite{ch2ssl}. A1 paper:
doi\string:10.5281/zenodo.23656804.

\paragraph{The lemma \lean{signed_endpoints} \tierAs.} It is the only lemma of the reduction proved completely. It
handles a Voronoi sector whose foot of the perpendicular falls outside its edge. Setting: a kernel
$G(d,s)$, even in $s$ and locally integrable for $d\ge d_0$, and its primitive $g_d(x)=\int_0^xG(d,s)\,ds$.

\begin{center}\small
\begin{tabular}{@{}p{0.34\linewidth}p{0.6\linewidth}@{}}\toprule
Manuscript & Lean (\lean{SignedEndpoints.lean:79-85}, verbatim)\\\midrule
$K<0$, $d_0>0$; $G(d,\cdot)$ even & \lstinline[style=chdeux]!(hK : K < 0) (hd₀ : 0 < d₀)!, \lstinline[style=chdeux]!heven : ∀ d, d₀ ≤ d → ∀ s, G d (-s) = G d s!\\
locally integrable & \lstinline[style=chdeux]!hint : ∀ d, d₀ ≤ d → ∀ a b : ℝ, IntervalIntegrable (G d) volume a b!\\
$2g_d(x)\ge K$ for $x\ge0$ & \lstinline[style=chdeux]!hpos : ∀ d, d₀ ≤ d → ∀ x, 0 ≤ x → K ≤ 2 * gPrim G d x!\\
$\int_0^\infty\max\{0,-G(d,s)\}\,ds\le-K$ & \lstinline[style=chdeux]!hneg : ∀ d, d₀ ≤ d → ∫⁻ s in Ioi 0, ENNReal.ofReal (-G d s) ≤ ENNReal.ofReal (-K)!\\
conclusion for $|\alpha|,|\beta|<\frac\pi2$, $\alpha+\beta>0$ & \lstinline[style=chdeux]!K ≤ gPrim G d (d * Real.tan α) + gPrim G d (d * Real.tan β)!\\\bottomrule
\end{tabular}
\end{center}
\begin{leanbox}[--- \lean{formal_cert/Programmes/SandierSerfaty/SignedEndpoints.lean:33, 79-85} (proof omitted)]
\begin{lstlisting}[style=chdeux]
noncomputable def gPrim (G : ℝ → ℝ → ℝ) (d x : ℝ) : ℝ := ∫ s in (0 : ℝ)..x, G d s
\end{lstlisting}
\begin{lstlisting}[style=chdeux]
theorem signed_endpoints (G : ℝ → ℝ → ℝ) (K d₀ : ℝ) (hK : K < 0) (hd₀ : 0 < d₀)
    (heven : ∀ d, d₀ ≤ d → ∀ s, G d (-s) = G d s)
    (hint : ∀ d, d₀ ≤ d → ∀ a b : ℝ, IntervalIntegrable (G d) volume a b)
    (hpos : ∀ d, d₀ ≤ d → ∀ x, 0 ≤ x → K ≤ 2 * gPrim G d x)
    (hneg : ∀ d, d₀ ≤ d → ∫⁻ s in Ioi 0, ENNReal.ofReal (-G d s) ≤ ENNReal.ofReal (-K)) :
    ∀ d, d₀ ≤ d → ∀ α β : ℝ, |α| < π / 2 → |β| < π / 2 → 0 < α + β →
      K ≤ gPrim G d (d * Real.tan α) + gPrim G d (d * Real.tan β) := by
\end{lstlisting}
\end{leanbox}
\noindent The integral of the negative part is a \lean{lintegral} ($\int^-$, with values in $[0,+\infty]$): an
infinite integral \emph{violates} the hypothesis, instead of producing the junk value $0$ of the Bochner integral.

\begin{theorem}[signed endpoints]\label{ch2:thm:signed}
Under the hypotheses of the table, for $d\ge d_0$ and $|\alpha|,|\beta|<\frac\pi2$ with $\alpha+\beta>0$:
$g_d(d\tan\alpha)+g_d(d\tan\beta)\ge K$. \tierAs
\end{theorem}
\begin{humanproof}
\emph{(1) Oddness.} By evenness of $G(d,\cdot)$ and the change of variable $s\mapsto-s$, $g_d(-x)=-g_d(x)$.

\emph{(2) Two nonnegative angles.} Then $x=d\tan\alpha\ge0$ and $y=d\tan\beta\ge0$; the hypothesis $2g_d\ge K$ gives
$g_d(x)\ge K/2$ and $g_d(y)\ge K/2$; add them.

\emph{(3) Opposite signs, $\alpha<0$.} Then $\alpha+\beta>0$ forces $\beta>-\alpha>0$. Set $u=-d\tan\alpha>0$ and
$v=d\tan\beta$; since $\tan$ is strictly increasing on $(-\frac\pi2,\frac\pi2)$, we have $u<v$. By oddness,
$g_d(d\tan\alpha)+g_d(d\tan\beta)=g_d(v)-g_d(u)=\int_u^vG(d,s)\,ds$.

\emph{(4) Central estimate.} $\int_u^vG\ge-\int_u^vG^-\ge-\int_0^\infty G^-\ge K$, where $G^-=\max(0,-G)$, by
the hypothesis on the negative part and $0<u$.

\emph{(5)} The case $\beta<0$ is symmetric; both negative is excluded by $\alpha+\beta>0$.
\end{humanproof}
\noindent Without $\alpha+\beta>0$, the statement is false (Exercise~\ref{ch2:ex:signed}); the refutation is checked by the
kernel (\lean{ctl_signed_without_sum_false}), and the same integral kernel $G\equiv1$ satisfies all the other hypotheses
(\lean{ctl_signed_hyps}), so that the theorem is not vacuous.

\begin{warnbox}[: what \lean{signed_endpoints} does not do]
The lemma is stated for an \emph{abstract} $G$. The concrete $G$ requires the radial data $f_0,e$ of Section~4 of the
manuscript, which are not defined in Lean, and nothing derives its two hypotheses from the certificate: the lemma \textbf{is
not connected to the certificate}. The preregistered stopping criterion applies: \textbf{no formal proof of the conjecture
is claimed}. The faithfulness of the statement to the manuscript was compared by an agent, not by a human.
\end{warnbox}

\subsection{B3: why the method does not carry over as is to dimension three}\label{ch2:ssec:b3}
In dimension three, the analogue is the Wigner conjecture: the body-centred cubic (bcc) lattice would minimise the
jellium energy among \emph{all} configurations; it is known to be optimal only among Bravais lattices of covolume one, for the Coulomb Born energy (Muratov--Novaga, Theorem~1, p.~4~\cite{ch2mn})
\tierL. The B3 study (doi\string:10.5281/zenodo.23655002)~\cite{ch2b3} splits the planar proof into 13 steps and classifies each one:
dimension-independent, transferable, or specifically planar. The incidence count (Euler's identity per
cell) is \emph{not} the obstacle.

\paragraph{The obstacle: exact calibration.} The whole planar proof is a lower bound via a trial potential $H$:
$W(h)=F(H)+\frac1{8\pi}\int|\nabla(h-H)|^2\ge F(H)$ (dual identity, \tierC), and the bound is optimal only if
$H$ \emph{equals} the true Green's function $h$ at the optimal lattice. In dimension two this is possible: in each sector,
$\cos6m\varphi=T_{6m}(d/t)$ (Chebyshev polynomial, $t=d/\cos\varphi$), so $h_\triangle$ depends on only two
variables and belongs to the trial class. In dimension three, the restriction of the bcc Green's function to a circle
of fixed radius on a face carries a $\cos6\psi$ mode (hexagonal faces) or a $\cos4\psi$ mode (square faces) \tierX, whereas
an axial $H$ is constant there \tierC.

\paragraph{Result, with its exact scope.} For the radial, axial and ``intrinsic hierarchical'' trial classes,
a projection argument \tierC, evaluated in floating point, forces a deficit: the certifiable constant is at most
$W_{\rm bcc}-\delta$ with $\delta\ge8.751\times10^{-5}$ (hierarchical) or $\delta\ge1.058\times10^{-4}$ (axial,
$1.18$ times the fcc--bcc gap $8.969\times10^{-5}$); these values are \tierX. For richer classes (which read
data from neighbouring faces), the negative result rests only on a counting heuristic (about 12 conditions for 3
multipliers, against 3 for 3 in dimension two) \tierC{} and \textbf{is not a theorem}. Nothing is proved about
bcc.

% ---------------------------------------------------------------------
\section{Exercises}\label{ch2:sec:ex}

Exercises marked \textbf{(Lean)} point to a kernel-checked reference solution in the repository
\tierA{}; the solutions sketched here are \tierC.

\begin{exercise}[dipolar control]\label{ch2:ex:hd}
Let $d>0$ and $h_d(t)=t^{-1/2}-2(t+d^2)^{-1/2}$. Compute $h_d(d^2)$ and show that it is strictly negative.
Deduce that $h_d$ is not completely monotone. Interpret physically: which pair-counting error would
produce $h_d$?
\end{exercise}

\begin{exercise}[a signed mixture]\label{ch2:ex:signedmix}
Let $\sigma(t)=e^{-t}-2e^{-2t}$. Show that $\sigma(\ln\frac32)=-\frac29$. Why does Theorem~\ref{ch2:thm:laplace}
require $m\ge0$? Exhibit a signed measure whose Laplace transform is $\sigma$.
\end{exercise}

\begin{exercise}[the tail bound]\label{ch2:ex:diffle}
Prove Lemma~\ref{ch2:lem:diffle} again using the mean value theorem applied to $x\mapsto x^{-1/2}$ on $[t,t+c]$.
Which constant does one obtain? Compare with $\frac c2$.
\end{exercise}

\begin{exercise}[the heart of the Yukawa tail bound]
(a) Show that $e^x\ge x^2/2$ for $x\ge0$. (b) Deduce $g_\kappa(t)\le\frac2{\kappa^2}t^{-3/2}$. (c) Show that,
in the lattice $\Lat$, the number of points in the annulus $k\le|a|<k+1$ is $O(k)$ (hint: Voronoi cells
of area $1$, disjoint, contained in an enlarged annulus), and conclude that $\sum_{a\ne0}|a|^{-3}<\infty$.
(d) Why does the argument fail for $\kappa=0$?
\end{exercise}

\begin{exercise}[moments of the Keldysh mixture]
Show that for $u\ge0$, $c>0$ and $N\in\N$, $u^Ne^{-cu}\le N!/c^N$. Deduce that
$a\mapsto(4a)^{-n}e^{-c/(4a)}\,\frac{w(a)}{2a}$ is bounded on $(0,1]$. (Lean: \lean{pow_mul_exp_neg_le}.)
\end{exercise}

\begin{exercise}[the $\top\le\top$ trap]
Show that for a planar lattice, $\sum_{0<|a|\le R}|a|^{-1}$ grows like $R$. Deduce why one
\emph{expects} $\mathrm{latticeEnergy}(\mathcal K_{r_0})=+\infty$. Why does this not make the theorem
\lean{triangular_keldysh_optimal} false, but only vacuous?
\end{exercise}

\begin{exercise}[quadrupoles]\label{ch2:ex:quad}
Let $f$ be CM and $c\ge0$, $g(t)=f(t)-f(t+c)$ and $\Delta_ah(t)=h(t)-h(t+a)$. (a) Check that
$3f(t)-4f(t+c)+f(t+4c)=\Delta_cg(t)+\Delta_{2c}g(t)+\Delta_{3c}g(t)$. (b) Deduce, using item (v) of
Section~\ref{ch2:sec:lib}, that the interaction of two quadrupoles is CM. (c) Show that the field of a single quadrupole,
$t^{-1/2}-2(t+1)^{-1/2}+(t+4)^{-1/2}$, is negative at $t=3$, using $\sqrt3>1.7$ and $\sqrt7>2.6$.
\end{exercise}

\begin{exercise}[CM or Bernstein?]
Show that: (a) $1/t$ is CM and positive but is not a Bernstein function; (b) $e^{-1/t}$ is not CM (hint:
the sign of its derivative); (c) $t^\alpha$ is a Bernstein function for $0\le\alpha\le1$ but not for $\alpha=2$. Explain
why (b) shows that, in Theorem~\ref{ch2:thm:subord}, the inner function cannot be just any
positive function.
\end{exercise}

\begin{exercise}[the hypothesis $\alpha+\beta>0$]\label{ch2:ex:signed}
With $G\equiv1$, $K=-1$, $d_0=d=1$: check the four hypotheses of \lean{signed_endpoints}, then show that for
$\alpha=\beta=-\frac\pi4$ one has $g_1(\tan\alpha)+g_1(\tan\beta)=-2<K$. Which step of the proof of
Theorem~\ref{ch2:thm:signed} uses $\alpha+\beta>0$?
\end{exercise}

\begin{exercise}[interval arithmetic]
With $Q=100$, encode $a=[-1.25, 0.33]$ as $(-125,33)$ and $b=[0.50, 0.67]$ as $(50,67)$. (a) Compute the four
endpoint products and the outward-rounded product interval. (b) With Euclidean integer division
$\lfloor\cdot\rfloor$, show that $x\mapsto-\lfloor -x/Q\rfloor$ computes $\lceil x/Q\rceil$. (c) Why would rounding
to nearest be incorrect for a certificate?
\end{exercise}

\begin{exercise}[\textbf{(Lean)} a negative control]\sloppy
Starting from \lean{h_neg_at : h d (d ^ 2) < 0} (for $d>0$), prove in Lean \lstinline[style=chdeux]!¬ AdmissiblePotential (h d)!.
Hint: \lean{intro H}, then use the second component \lean{H.2.1} of the definition. Reference solution:\newline
\lean{formal_cert/Autoresearch/DipolarExciton.lean:113-116} \tierA.
\end{exercise}

\begin{exercise}[\textbf{(Lean)} corollaries of the library]
(a) State in Lean that $t\mapsto\mathrm{riesz}\,s\,t-\mathrm{riesz}\,s\,(t+c)$ is CM for $s,c\ge0$ and prove it in
one line using \lean{riesz_cm} and \lean{sub_comp_add_const}. Reference solution:
\lean{Corollaries.lean:33-35} \tierA. (b) Same question for $e^{-at^\alpha}$, $a\ge0$, $0\le\alpha\le1$, using
\lean{exp_neg_comp} and \lean{IsBernsteinOn.rpow}. Reference solution: \lean{Corollaries.lean:59-61} \tierA.
(c) Why is the condition \lstinline[style=chdeux]!(fun _ ht => Real.rpow_pos_of_pos ht α)! required?
\end{exercise}

\begin{exercise}[($\dagger$) University exercise: Laplace representations of Riesz and dipolar potentials]
(a) For $s>0$ and $t>0$, show that $t^{-s}=\frac1{\Gamma(s)}\int_0^\infty u^{s-1}e^{-tu}\,du$ (substitute $v=tu$ in
the integral defining $\Gamma(s)$). Deduce directly, by differentiation under the integral sign, that $t^{-s}$ is
completely monotone, and recover the formula $(-1)^n\frac{d^n}{dt^n}t^{-s}=s(s+1)\cdots(s+n-1)\,t^{-s-n}$.
(b) Using (a) with $s=\frac12$ and $\Gamma(\frac12)=\sqrt\pi$, prove the identity of the remark ``Bernstein reading'':
for $c\ge0$, $t^{-1/2}-(t+c)^{-1/2}=\pi^{-1/2}\int_0^\infty e^{-ut}\,(1-e^{-uc})\,u^{-1/2}\,du$, and deduce once more
that $\mathrm{shiftDiff}_{1,c}$ is CM.
(c) More generally, show that if $f(t)=\int_{[0,\infty)}e^{-tu}\,\mu(du)$ for a positive measure $\mu$, then for $c\ge0$
the function $f(t)-f(t+c)$ is the Laplace transform of the positive measure $(1-e^{-cu})\,\mu(du)$.
\end{exercise}

\begin{exercise}[($\dagger$) University exercise: Jacobi's theta identity and a square-lattice Gaussian sum]
Take for granted the Poisson summation formula on $\R$: for a Schwartz function $\phi$, $\sum_{n\in\Z}\phi(n)=\sum_{k\in\Z}\widehat\phi(k)$
with $\widehat\phi(\xi)=\int_\R\phi(x)e^{-2\pi i x\xi}\,dx$.
(a) Show that the Fourier transform of $x\mapsto e^{-\pi ax^2}$ ($a>0$) is $\xi\mapsto a^{-1/2}e^{-\pi\xi^2/a}$.
(b) Deduce Jacobi's identity $\theta(a)=a^{-1/2}\,\theta(1/a)$, where $\theta(a)=\sum_{n\in\Z}e^{-\pi an^2}$.
(c) Show that $\sum_{x\in\Z^2}e^{-\pi a\|x\|^2}=\theta(a)^2$, and hence
$\sum_{x\in\Z^2}e^{-\pi a\|x\|^2}=a^{-1}\sum_{x\in\Z^2}e^{-\pi\|x\|^2/a}$.
(d) Deduce that the Gaussian lattice energy of the square lattice, $E(a)=\sum_{x\in\Z^2\setminus\{0\}}e^{-\pi a\|x\|^2}$,
satisfies $E(a)=a^{-1}-1+O(a^{-1}e^{-\pi/a})$ as $a\to0^+$, and $E(a)\sim4e^{-\pi a}$ as $a\to+\infty$.
\emph{Answer.} d) $E(a)=a^{-1}\theta(1/a)^2-1$ and $\theta(1/a)^2=1+4e^{-\pi/a}+O(e^{-2\pi/a})$.
\end{exercise}

\label{ch2:fin}
